% Working draft for MNRAS.
% Canonical copy: https://www.overleaf.com/project/6abf692aab636771b039c71a
% The measurement is the stellar linking length matched to an isophote
% in group-mass haloes. The Jacobi-radius finder paper is the separate file ms_jacobi.tex.
% Author order: Jaehyun Lee, Jihye Shin, Yonghwi Kim, Jaewon Yoo,
% M. James Jee, Jongwan Ko, Juhan Kim (corresponding).
% The group-mass sample is finished. The supplement documents the
% additional TNG300-1 merger-tree growth analysis.
\documentclass[usenatbib]{mnras}
\usepackage{graphicx}
\usepackage{amsmath,amssymb}
\usepackage[T1]{fontenc}
\usepackage{newtxtext,newtxmath}
\usepackage{booktabs}
\usepackage{placeins}
\hyphenation{pixels}

\usepackage{xr}
\externaldocument[main-][nocite]{main}
\graphicspath{{figures/}}
\title{Supplementary methods and numerical checks}
\author{Lee et al.}
\begin{document}
\maketitle
\section{Simulation resolution and subgrid settings}
The RAMSES mesh spans levels 10 to 17. At $z=0.227$, the output length
unit is $53.393\,\mathrm{Mpc}$, the proper box size. The finest cell is
$2^{-17}$ of the output length unit, or $0.407$ proper kpc. Proper cell
size decreases at earlier outputs in proportion to the scale factor.
In periodic catalogue coordinates, the same cell spans
$44.171\times 2^{-17}=0.337\,h^{-1}\mathrm{ckpc}$ at every output.

Star particles form above a density threshold with a minimum birth mass
of $10^{5}\,\mathrm{M}_{\odot}$. The threshold is
$n_{\star}=0.1\,\mathrm{H\,cm}^{-3}$, the star-formation time-scale is
$t_{\star}=1.632\,\mathrm{Gyr}$, and the polytropic temperature is
$T_{2,\star}=15\,\mathrm{K}$. The polytropic index is $g_{\star}=1$,
so the temperature floor does not rise with density. Type II supernova
feedback is mechanical, with $E_{\mathrm{SNII}}=7\times10^{51}\,\mathrm{erg}$,
$\eta_{\mathrm{SN}}=0.20$, and a $5\,\mathrm{Myr}$ delay. Type Ia
feedback is thermal. Black-hole seeds of $5\times10^{4}\,\mathrm{M}_{\odot}$
are included, with spin evolution and magnetically arrested jet
switches on.

\section{Point-spread function transition}
\label{sec:psfwing}
The $A/r^{2}$ wing has zero amplitude inside $7\,\mathrm{arcsec}$,
rises linearly from $7$ to $8\,\mathrm{arcsec}$, and reaches full
amplitude beyond $8\,\mathrm{arcsec}$. The wing is added to the Moffat
core before the kernel is truncated at $80\,\mathrm{arcsec}$ and
renormalised to one. Main-paper Section~\ref{main-sec:psf} gives the
core width and band-dependent wing amplitudes.

\section{Galaxy membership before the stellar link}
The friends-of-friends group in main-paper Section~\ref{main-sec:group} is grown from a
galaxy already defined by PGalF \citep{kim2023}.
Figure~\ref{fig:finder} shows the galaxy definition on a schematic stellar
density field.
A density core is a local maximum surviving the watershed.
From each maximum the finder walks downhill until the basin touches
another maximum.
The density at the contact is the saddle.
The watershed is the boundary passing through the saddle.
A peak with prominence
$(\rho_{\mathrm{peak}}-\rho_{\mathrm{saddle}})/\rho_{\mathrm{peak}}$
below $0.30$ is folded into the higher peak touched first.
A basin with fewer than $100$ particles is discarded.
The host is the core with the highest stellar peak density.

Membership then walks outward through iso-density shells.
A shell is the set of particles lying between two successive density
thresholds, and the walk starts at the higher threshold.
When a shell encloses a satellite core, the tidal radius is
\begin{equation}
r_{t}
=
d\left(\frac{m(<r_{t})}{3\,M(d)}\right)^{1/3}.
\label{eq:rt}
\end{equation}
The radius is capped at $d/2$.
The separation of the host peak and the satellite peak is $d$.
The satellite mass $m(r)$ is the mass of stars, dark matter, and gas
around the satellite peak and inside radius $r$.
Particles already belonging to another core are omitted from $m(r)$.
The host mass $M(d)$ is the mass inside distance $d$ of the host peak
after removal of the particles already assigned to the satellite.
$M(d)$ sets the tidal boundary.
The satellite potential uses only $m$.
Particles outside $r_{t}$ return to the unassigned pool.
The host tests the unassigned pool in a later pass inside a search radius
up to the virial radius of the friends-of-friends halo.
The nucleus is the set of stars already belonging to the satellite
and lying within $0.012\,h^{-1}\mathrm{cMpc}$ of the satellite core.
At least five nucleus stars are required, and the nucleus stars set the
satellite rest frame.
The nucleus drawn in Figure~\ref{fig:finder} is enlarged, and the
circle remains visible.
The linking length in main-paper Section~\ref{main-sec:group} is applied after the
membership above.
The linking length is a separate stellar friends-of-friends boundary.

\begin{figure}
\centering
\includegraphics[width=\columnwidth]{figures/finder_geometry.pdf}
\caption{Boundaries defining a galaxy before the stellar linking
length is applied.
The density is a schematic with two surviving peaks and one weak peak,
drawn for a clear separation of the boundaries.
Panel (a) is a slice through the watershed.
The finder walks downhill from each local maximum until the basin meets
another maximum, and the density at the contact is the saddle.
The dashed curve is the watershed.
A surviving maximum is a density core.
The weak peak is folded into the host because the weak-peak prominence
$(\rho_{\mathrm{peak}}-\rho_{\mathrm{saddle}})/\rho_{\mathrm{peak}}$
is below $0.30$.
On the drawing the weak-peak prominence is $0.18$.
The host is the core with the higher peak density.
The small circle on the satellite is the nucleus setting the satellite rest
frame, drawn larger than the true radius of
$0.012\,h^{-1}\mathrm{cMpc}$.
Panel (b) is the same slice during membership.
The solid contours are three density thresholds, with
$\rho_{\mathrm{high}}>\rho_{\mathrm{mid}}>\rho_{\mathrm{low}}$.
A shell is the set of particles between two successive thresholds, and
the arrow runs from the higher threshold toward the lower one.
The circle centred on the satellite is the tidal radius $r_{t}$ of
Equation~\eqref{eq:rt}.
On the drawing the tidal radius lies inside the cap at $d/2$.
Filled symbols inside $r_{t}$ can be kept with the satellite.
Open symbols outside $r_{t}$ return to the unassigned pool.
Panel (c) shows the same cores in three dimensions.
The small closed curves are the high-density shells around each core,
the outer curve is the lower threshold, the sheet between the cores is
the watershed, and the sphere is the tidal radius.
The linking length of main-paper Section~\ref{main-sec:group} is applied after the
watershed membership and is a separate stellar friends-of-friends boundary.}
\label{fig:finder}
\end{figure}

\section{Linking distance and density}
For equal-mass stars the estimate $3\bar\rho/(2\pi b_{\mathrm{link}}^{3})$ is about
$60\,\bar\rho$ at $b_{\mathrm{link}}=0.2$.
The estimate is an equal-mass analogy.
The estimate provides neither a calibrated density boundary for the
unequal-mass stars nor a single isodensity contour.
The percolation critical density is instead $81.62\,\bar\rho$ at the same
linking length, and the boundary is a range around the critical density
\citep{more2011}.
The critical density is not rescaled to any other linking length.

\section{Sample definitions and transfer checks}
The full Darwin catalogue contains 76 halo-output entries at outputs
00081, 00178, 00475, 00786, 00967, and 01001. We retain 75 entries for
calibration because one $z=0.227$ entry has a maximum at a tested-grid
endpoint in both-band selection. The artificial-redshift experiment
instead uses every entry at $z\leq1.990$. The experiment omits the one
$z=3.002$ entry and includes the endpoint entry at $z=0.227$. The
complete IDs, masses, and inclusion flags are listed in
\texttt{sample\_manifest.tsv}. The calibration and artificial-redshift samples are distinct.

Exact dark-matter particle IDs establish links between selected haloes
through output 00786 ($z=0.500$). Output 00967 has a top-100 FoF
member-particle file, but the candidate 32-bit ID field fails the identity
checks. A stratified sample reads three 4,096-record segments from each of
the ten selected groups. The candidate ID is zero in 15.2--59.7 per cent of
the 12,288 sampled records per group, and every sample contains repeated
values. The output 01001 member-particle file is 68 bytes long and contains
no usable particle payload. The nine apparent matches between outputs 00967
and 01001 therefore use positions and remain provisional. The catalogues do
not provide a merger tree. We report output entries separately and do not
treat every entry as an independent system. For the exact-ID check, an
earlier--later pair is linked when the later group is the best-overlap
descendant and contains more than half of the earlier group's dark-matter
particles. The 20 accepted links form 33 connected groups among 53 entries
in outputs 00178, 00475, and 00786. Entries without an accepted link remain
separate. The group count is not a complete count of independent lineages.

The leave-one-output test uses the corrected-distance photometry
adopted in the main paper. In each test, we exclude one output and
catalogue centres within $0.5\,h^{-1}\mathrm{cMpc}$, then choose one
$b_{\mathrm{link}}$ from the remaining entries. All six tests select
$b_{\mathrm{link}}=0.10$ in both bands. Across the 75 excluded entries,
median overlap is $0.9582$ in $V$ and $0.9626$ in $g$. The median gain
over $b_{\mathrm{link}}=0.20$ is $0.0363$ and $0.0339$. This test
measures transfer between outputs from one volume, not between
independent haloes. Results for each test are in
\texttt{validation\_corrected\_predictions.tsv}.

For a second test, we omit all entries in one exact-ID connected group
at a time and give each of the remaining groups equal weight when
choosing $b_{\mathrm{link}}$. We also exclude any other group with a
member within $0.5\,h^{-1}\mathrm{cMpc}$ of an omitted member. All 33
tests select $b_{\mathrm{link}}=0.10$ in both bands. The median overlap
for the omitted groups is $0.9556$ in $V$ and $0.9621$ in $g$. The equal-group mean gain over
$b_{\mathrm{link}}=0.20$ is $0.0487$ ($0.0415$--$0.0555$) in $V$ and $0.0415$
($0.0336$--$0.0494$) in $g$; intervals are 16th--84th percentiles from
5,000 group-bootstrap draws. This result is restricted to the three
outputs with exact particle-ID information.

A sensitivity check adds reciprocal-nearest positional matches for
intervals without exact particle IDs. The 75 retained entries form 35
provisional groups; all 35 comparisons again select $b_{\mathrm{link}}=0.10$
in both bands. The equal-group mean gain over $b_{\mathrm{link}}=0.20$ is $0.0423$
($0.0359$--$0.0486$) in $V$ and $0.0379$ ($0.0308$--$0.0448$) in $g$.
Because some links use positions rather than particle IDs, these values
are conditional on the provisional grouping and do not establish
independent-halo generalisation.

\begin{table*}
\centering
\caption{Image centres, isophotal detection rules, and treatment of failed
contours in the simulation samples. A central-pixel detection means that the
connected above-threshold region contains the image-centre pixel.}
\label{supp:centredetection}
\small
\begin{tabular}{p{2.5cm}p{3.5cm}p{5.0cm}p{5.0cm}}
\toprule
Sample & Image centre & Detection rule & Failure treatment \\
\midrule
Darwin calibration, 75 entries &
Mass-weighted centre of the dark matter, gas, sink particles, and stars in
the catalogue central galaxy &
Use the connected threshold region at the centre pixel; if that pixel is
fainter than the limit, use the region around the nearest pixel above the
limit. An empty image gives an empty region. &
Nearest-pixel selections are retained and flagged. An entry is excluded from
calibration if any view has its maximum at either end of the tested linking
grid. \\
CAMELS CV20--CV26, 49 hosts &
Minimum-potential position of the central subhalo &
Use the same threshold and connected-region rule as for Darwin. &
No target contour failed in the reported fixed-parameter image checks. \\
CAMELS CV0--CV19, 130 hosts &
Minimum-potential position of the central subhalo &
The central component must be detected; a disconnected nearest-bright
component is not accepted as the central contour. &
Fallback, empty or zero-luminosity regions, and regions touching the image
edge are invalid. Their IoU and completeness are set to zero; luminosity and
$R_{50}$ errors are undefined. \\
TNG300-1 transfer, 40 hosts &
\texttt{SubhaloPos}, the position of the maximally bound particle &
Use the connected threshold region at the centre; a missing region is not a
detection. &
Missing regions and regions touching the image edge are invalid for the
transfer summary. \\
TNG300-1 growth, 36 branches &
\texttt{SubhaloPos} at each output &
Nearest-pixel selections are retained in the full sample and flagged. The
detection-qualified sample requires a central-pixel detection in all views
at both endpoint outputs. &
The full sample retains fallback selections. The endpoint-detection cut
leaves 36 branches in $V$ and 29 in $g$. \\
\bottomrule
\end{tabular}
\end{table*}
\FloatBarrier

The image centres are defined by simulation catalogues and particle
distributions, not recovered from the mock images. The validation samples
therefore test isophotal selection at known centres; they do not test an
image-based galaxy-centre finder. The main text and
Table~\ref{supp:centredetection} distinguish central-pixel detection from
other invalid-contour rules.

\subsection{Corrected-distance fine grid and PSF-aperture weights}
\label{supp:psfweights}

We compare two light-overlap measures using the same corrected-distance
isophotal regions and the same stellar-group memberships. In the positional
measure, a star contributes its full flux when its projected position is
inside the isophotal region. In the PSF-aperture measure, its flux is
multiplied by the fraction of its normalized PSF kernel that falls inside
that region. This second measure estimates the stellar flux inside the
aperture after PSF spreading. It changes the overlap weights, not group
membership. We evaluate both measures on the 15-point linking-parameter
grid and choose a best value separately for each projection. We first take
the median of the three projection-specific best values for each halo
entry, then report the distribution across the 75 entries.

The median fine-grid best parameter is $0.09$ in both bands for both
weighting rules. Table~\ref{supp:psf-grid} gives the percentile ranges and
overlap scores. For the median-over-entry score curves, the positional
measure peaks at $b_{\mathrm{link}}=0.09$ in both bands; the PSF-aperture
measure peaks at $0.09$ in $V$ and $0.10$ in $g$. The PSF-aperture scores
are lower because flux from stars outside the isophote is spread into its
aperture and flux from stars inside it is spread out. The preferred
linking lengths change little, and $b_{\mathrm{link}}=0.10$ outperforms
$0.20$ for both weighting rules.

Applying the central-pixel detection rule in all three views retains 55
entries in $V$ and 53 in $g$. Their median coarse-grid best parameters
remain $0.08$ and $0.10$, respectively, while both fine-grid weighting
rules give a median of $0.09$ in both bands. At the entries' original
coarse-grid selected parameters, the median PSF-aperture IoU is $0.950$
in $V$ and $0.949$ in $g$ for this subset, compared with $0.944$ and
$0.950$ for all 75 entries. The complete central-detection flags and
per-entry scores are in
\texttt{supplement/astra\_review\_response\_20261106/results/}.

\begin{table*}
\centering
\caption{Corrected-distance fine-grid overlap under positional and
PSF-aperture flux weighting. For each halo entry, the three views are
reduced before percentiles are calculated across entries. The listed
best-parameter range is the 16th--84th percentile; the optimized overlap
is the median at each view's individual maximum.}
\label{supp:psf-grid}
\small
\begin{tabular}{llrrrr}
\toprule
Band & Weighting & Median $b_{\mathrm{link}}$ (16th--84th) &
Median optimized IoU & IoU at $b_{\mathrm{link}}=0.10$ & IoU at $b_{\mathrm{link}}=0.20$ \\
\midrule
$V$ & Position & $0.09\;(0.07$--$0.12)$ & 0.968 & 0.958 & 0.920 \\
$V$ & PSF aperture & $0.09\;(0.07$--$0.11)$ & 0.946 & 0.937 & 0.905 \\
$g$ & Position & $0.09\;(0.07$--$0.12)$ & 0.973 & 0.963 & 0.933 \\
$g$ & PSF aperture & $0.09\;(0.07$--$0.11)$ & 0.951 & 0.944 & 0.917 \\
\bottomrule
\end{tabular}
\end{table*}

\begin{figure*}
\centering
\includegraphics[width=\textwidth]{supplement/astra_review_response_20261106/results/corrected_positional_and_psf_overlap.pdf}
\caption{Corrected-distance median light-weighted overlap on the 15-point
linking-parameter grid. The blue curves assign each star's full flux by
projected position; the orange curves weight its flux by the PSF fraction
inside the isophotal aperture. Curves show the median of the three-view
median for each halo entry; shaded regions show the 16th--84th percentiles
across entries.}
\label{fig:supp-psf-grid}
\end{figure*}

\subsection{Within-output validation at $z=0.1904$}
For the within-output test, we omit one of the 12 haloes at $z=0.1904$
and choose a common parameter from the other 11. Each calibration halo
contributes the median of its three projection-specific parameters with
maximum IoU; we then take the median across the 11 haloes. All 12 tests
select $b_{\mathrm{link}}=0.12$ in $V$ and $0.10$ in $g$. The results
for each omitted halo and projection are in
\texttt{supplement/within\_output\_validation\_z019\_20261106/}.
Because the tests share haloes and use one output, they do not provide
independent-volume uncertainty.

The output-weighting sensitivity computes the median-over-halo IoU curve
within each of the six outputs, then averages the six curves with equal
weights. The best tested parameter is $b_{\mathrm{link}}=0.10$ in both
bands. The output-specific measurements and plotting script are in
\texttt{supplement/output\_weighting\_sensitivity\_20261106/}.

\section{CAMELS Swift-EAGLE transfer sample}
The fixed-$b_{\mathrm{link}}$ transfer test contains 49 host groups from first-generation
CAMELS Swift-EAGLE L25n256 CV boxes 20--26, using snapshot 090 and the
associated \texttt{groups\_090.hdf5} catalogue. The 25 $h^{-1}$ cMpc box
contains $256^3$ initial gas elements and dark-matter particles. Snapshot
header mass units give an initial gas-element mass of
$1.89\times10^7\,\mathrm{M}_{\odot}$. Star particle masses evolve after stellar
mass loss. The catalogue records a maximum proper stellar softening of
$5.84\,h^{-1}$ pkpc. The cosmology is $\Omega_m=0.3$, $h=0.6711$.
For every host and band, the same simulation potential minimum centres
the image, the 30 pkpc aperture, and the projected radial measurements.
The aperture is a two-dimensional circle with projected radius at most
30 pkpc. For each selected stellar set, projected $R_{50}$ is the
smallest projected particle radius enclosing at least half of the total
stellar flux of the selected set. Luminosity and radius errors are respectively
$|\log_{10}(L_{\rm selected}/L_{\rm isophote})|$ and
$|\log_{10}(R_{50,\rm selected}/R_{50,\rm isophote})|$. We use the
median over the three projections for each host, then report the
90th percentile across hosts. The PSF-convolved image defines the
isophotal region. Full particle fluxes at projected positions define
luminosity and $R_{50}$, without PSF-aperture weights. The transfer test
compares one finite-resolution simulation and does not demonstrate
convergence with particle or force resolution.
The simulation centre is known, so the test does not evaluate automatic
centre finding. Satellite galaxies are masked consistently. Full host
metrics, per-box summaries, and paired host contrasts are supplied as
\texttt{camels\_host\_metrics.csv}, \texttt{camels\_box\_summary.csv},
and \texttt{camels\_paired\_contrasts.csv}.

For $b_{\mathrm{link}}=0.10$, the mean host-median overlaps are 0.894 in $V$ and 0.872
in $g$. The corresponding host medians are 0.893 and 0.872. Across the
seven CV boxes, mean host overlap ranges from 0.872 to 0.904 in $V$
and 0.842 to 0.886 in $g$. Relative
to $b_{\mathrm{link}}=0.20$, paired mean improvements are 0.094 (95 per cent
box-then-host bootstrap interval 0.083--0.106) and 0.111 (0.098--0.127).
We first reduce the projections to one statistic per host. The bootstrap
resamples CV boxes and then hosts within each sampled box. At $b_{\mathrm{link}}=0.10$, the
90th-percentile absolute errors are 0.033 and 0.051 dex in luminosity,
and 0.048 and 0.081 dex in projected half-light radius. The 30 pkpc
aperture gives errors of 0.166/0.158 dex in luminosity and 0.238/0.236 dex
in radius ($V/g$). The fixed linking parameter reduces luminosity and
radius errors for the sample with known centres, although the stellar
selections still differ for individual hosts.

We also compare $b_{\mathrm{link}}=0.10$ directly with the 30 pkpc aperture on the same
49 hosts. The mean paired overlap difference is $+0.040$ in $V$
(95 per cent box-then-host interval $+0.005$--$+0.083$) and $+0.010$
in $g$ ($-0.024$--$+0.051$). The mean selected-light fraction outside
the isophote increases by $0.055$ and $0.074$, while completeness
increases by $0.094$ and $0.081$. The linking selection recovers more
isophotal light and gives smaller luminosity and radius errors, but also
includes more light outside the isophote. The paired aperture table
\texttt{camels\_b010\_vs\_aperture.csv} contains paired differences in
signed biases and absolute errors, with 5,000-draw host and box-then-host
intervals.

\section{Test in separate CV0--CV19 boxes}

We selected 130 FoF hosts from the first-generation CAMELS Swift-EAGLE
L25n256 boxes CV0--CV19 at snapshot 090. The FoF masses span
$10^{13}$--$10^{14}\,h^{-1}\mathrm{M}_{\odot}$, and each central subhalo
contains at least 100 star particles. The boxes do not overlap the earlier
CV20--CV26 sample. We rendered the $z=0$ systems at
$z_{\rm obs}=0.19044$ with the same filters, brightness limits, image scale,
PSF, satellite masks, and particle-flux calculation as the earlier CAMELS
analysis. Catalogue particle IDs define the Swift-EAGLE FoF membership.
The central subhalo's potential minimum, \texttt{SubhaloPos}, centres each
image, aperture, and projected radius. We compare $b_{\mathrm{link}}=0.20$,
common fixed $b_{\mathrm{link}}=0.10$, the low-redshift Darwin values
($b_{\mathrm{link}}=0.12$ in $V$ and 0.10 in $g$), and a projected 30 pkpc
aperture. We do not tune the linking parameter on these 130 hosts.

We did not use the 130-host measurements to choose a linking parameter.
The centre rule came from earlier CV20--CV26 and TNG300-1 work, so this
is not a test whose full analysis plan was set in advance. It measures transfer across boxes in
one Swift-EAGLE model, not across simulation models. Each host contributes
the median across three orthogonal projections. Bootstrap intervals use 2,000 draws, resampling boxes first and hosts
within sampled boxes second. A large failure is a host with an invalid target
contour in any view, median contamination above 0.20, or a factor-of-two
error in host-median luminosity or $R_{50}$.

\begin{table*}[t]
\centering
\footnotesize
\begingroup
\setlength{\tabcolsep}{3pt}
\begin{tabular*}{\textwidth}{@{\extracolsep{\fill}}llcccccc@{}}
\toprule
Band & Method & Mean IoU (95\% interval) & Completeness & Contamination &
\shortstack{$P_{90}(|\Delta\log L|)$\\(dex)} &
\shortstack{$P_{90}(|\Delta\log R_{50}|)$\\(dex)} & Large failures \\
\midrule
$V$ & 30 pkpc aperture & 0.837 [0.812, 0.862] & 0.846 & 0.009 & 0.180 & 0.252 & 5/130 \\
$V$ & $b_{\mathrm{link}}=0.20$ & 0.783 [0.767, 0.798] & 0.988 & 0.209 & 0.135 & 0.247 & 70/130 \\
$V$ & Darwin $b_{\mathrm{link}}=0.12$ & 0.864 [0.853, 0.875] & 0.966 & 0.107 & 0.063 & 0.094 & 5/130 \\
$V$ & Fixed $b_{\mathrm{link}}=0.10$ & 0.884 [0.874, 0.893] & 0.948 & 0.070 & 0.038 & 0.061 & 0/130 \\
$g$ & 30 pkpc aperture & 0.844 [0.820, 0.868] & 0.859 & 0.016 & 0.160 & 0.231 & 4/130 \\
$g$ & $b_{\mathrm{link}}=0.20$ & 0.742 [0.724, 0.759] & 0.989 & 0.251 & 0.179 & 0.303 & 99/130 \\
$g$ & Darwin $b_{\mathrm{link}}=0.10$ & 0.861 [0.849, 0.872] & 0.950 & 0.097 & 0.055 & 0.082 & 3/130 \\
\bottomrule
\end{tabular*}
\endgroup
\caption{CV0--CV19 validation sample metrics. Each host contributes the
median across three projections. Mean IoU intervals use the nested box-host
bootstrap. Completeness and contamination are equal-host means. The final
column counts hosts meeting the large-failure definition in the text. The
$g$-band Darwin calibration is $b_{\mathrm{link}}=0.10$ and is listed once.}
\label{tab:cv019-validation}
\end{table*}

All 130 catalogue particle-ID lists resolve to snapshot stars. The median
Jaccard similarity between catalogue membership and snapshot
\texttt{FOFGroupIDs} is 0.998. Five hosts in CV8 or CV11 have no overlap
between the two membership representations. The primary calculation retains
the catalogue particle IDs. Removing the five hosts changes mean IoU by at
most 0.0026, mean contamination by at most 0.0007, the 90th-percentile
luminosity error by at most 0.007 dex, and the large-failure fraction by at
most 0.006. The checks confirm that the catalogue IDs resolve to snapshot
stars, but do not show whether the disagreement comes from halo-index ordering
or a difference in the FoF assignments. The cause remains unidentified.

Supplementary tables give host, projection, box, and paired-comparison
measurements, the sample selection and membership checks, and the
analysis scripts in
\url{supplement/cv0_19_holdout_20261103/}. The scripts require the
original CAMELS snapshot and catalogue products.

\section{TNG300-1 SubhaloPos-centred transfer test}

We selected 40 FoF hosts from snapshot 84 of TNG300-1 ($z=0.1973$)
\citep{nelson2019tng,pillepich2018tng,springel2018tng,nelson2018tng,naiman2018tng,marinacci2018tng}.
FoF masses span $3.09\times10^{13}$--$1.288\times10^{14}\,
h^{-1}\mathrm{M}_{\odot}$, with a median of
$4.680\times10^{13}\,h^{-1}\mathrm{M}_{\odot}$. For each host, we select the
qualifying subhalo with the largest stellar-particle count. We centre images
and apertures on \texttt{SubhaloPos}, the position of the maximally bound
particle. TNG guidance recommends \texttt{SubhaloPos} for galaxy centring.
The all-type mass-weighted centre \texttt{SubhaloCM} can be displaced by
asymmetric or distant material.\footnote{The IllustrisTNG data specifications
define \texttt{SubhaloPos} and \texttt{SubhaloCM} at
\url{https://www.tng-project.org/data/docs/specifications/}. The position
guidance appears at
\url{https://www.tng-project.org/data/forum/topic/295/how-to-get-galaxy-positions/}.}
The stellar-density peak remains the root for linked-group selection. The
same satellite mask, filters, thresholds, image scale, PSF, and flux
calculation apply to Darwin and TNG. We compare the low-redshift Darwin
calibration ($b_{\mathrm{link}}=0.12$ in $V$ and 0.10 in $g$), fixed values
of 0.10 and 0.20, and a projected 30 pkpc aperture. The leave-one-host-out
test selects the same values as the low-redshift calibration, so we report
those transfer measurements once. Each host contributes the
median across three projections. Host haloes are the statistical units.
Earlier TNG overlap measurements using a different image-centre rule had already been
inspected. The present centre comparison is post hoc, and no linking
parameter was tuned using TNG measurements.

Table~\ref{tab:tng300-transfer} reports median IoU, median contamination,
mean signed luminosity and $R_{50}$ errors, 90th-percentile absolute errors,
and large-failure counts. Signed errors are
$\Delta\log L=\log_{10}(L_{\rm selected}/L_{\rm isophote})$ and
$\Delta\log R_{50}=\log_{10}(R_{50,\rm selected}/R_{50,\rm isophote})$.
The projected $R_{50}$ encloses half of the selected particle flux.
Contamination is one minus light-weighted precision. A large failure is a
factor-of-two error in luminosity or $R_{50}$, contamination above 0.20, or
a missing or edge-touching selected contour.

\begin{table*}[t]
\centering
\footnotesize
\begingroup
\setlength{\tabcolsep}{2.5pt}
\begin{tabular*}{\textwidth}{@{\extracolsep{\fill}}llccccccc@{}}
\toprule
Band & Method & Median IoU & Mean $\Delta\log L$ &
$P_{90}(|\Delta\log L|)$ & Mean $\Delta\log R_{50}$ &
$P_{90}(|\Delta\log R_{50}|)$ & Median contamination & Failures \\
\midrule
$V$ & $b_{\mathrm{link}}=0.20$ & 0.833 & +0.074 & 0.103 & +0.128 & 0.194 & 0.158 & 9/40 \\
$V$ & Fixed $b_{\mathrm{link}}=0.10$ & 0.893 & -0.027 & 0.059 & -0.037 & 0.078 & 0.025 & 0/40 \\
$V$ & Darwin low-$z$ & 0.892 & +0.002 & 0.031 & +0.009 & 0.057 & 0.056 & 0/40 \\
$V$ & 30 pkpc aperture & 0.716 & -0.153 & 0.245 & -0.221 & 0.334 & 0.000 & 7/40 \\
$g$ & $b_{\mathrm{link}}=0.20$ & 0.783 & +0.100 & 0.142 & +0.162 & 0.233 & 0.209 & 22/40 \\
$g$ & Fixed $b_{\mathrm{link}}=0.10$ & 0.883 & -0.007 & 0.042 & -0.003 & 0.065 & 0.053 & 0/40 \\
$g$ & Darwin low-$z$ & 0.883 & -0.007 & 0.042 & -0.003 & 0.065 & 0.053 & 0/40 \\
$g$ & 30 pkpc aperture & 0.737 & -0.139 & 0.243 & -0.188 & 0.304 & 0.000 & 5/40 \\
\bottomrule
\end{tabular*}
\endgroup
\caption{Host-level transfer metrics for the 40 TNG300-1 haloes. Each host
is represented by the median across three projections. The low-redshift
Darwin calibration uses $b_{\mathrm{link}}=0.12$ in $V$ and 0.10 in $g$;
the leave-one-host-out test selects these same values. The aperture is a
projected circle of radius 30 pkpc.}
\label{tab:tng300-transfer}
\end{table*}

The primary Darwin values increase median IoU over $b_{\mathrm{link}}=0.20$
by 0.067 in $V$ (95 per cent host-bootstrap interval 0.055--0.073) and
0.104 in $g$ (0.093--0.109). Median contamination falls by 0.093 in $V$
and 0.143 in $g$. Median absolute luminosity errors fall by 0.054 and
0.073 dex. Median absolute $R_{50}$ errors fall by 0.098 and 0.124 dex.
Large failures fall by 22.5 percentage points in $V$ and 55 points in $g$.
The linking selection has higher IoU and smaller luminosity and $R_{50}$
errors than the 30 pkpc aperture. The aperture selects less light outside
the isophote. The calibrated linking method therefore trades higher
contamination for closer agreement with the full isophotal light.

The proper distance between \texttt{SubhaloCM} and \texttt{SubhaloPos} has
a median of 72.4 pkpc, a 90th percentile of 171.9 pkpc, and a maximum of
396.7 pkpc. The stellar-density peak lies closer to \texttt{SubhaloPos},
with median, 90th-percentile, and maximum offsets of 0.62, 1.44, and
9.83 pkpc. At fixed $b_{\mathrm{link}}=0.10$, the earlier
\texttt{SubhaloCM}-centred rendering has 90th-percentile absolute
luminosity errors of 3.36 dex in $V$ and 3.30 dex in $g$. The present
\texttt{SubhaloPos}-centred rendering gives 0.059 and 0.042 dex. The earlier
stellar-density-peak-centred and current \texttt{SubhaloPos}-centred
renderings agree within 0.002 in host-level median IoU for 90 per cent of
hosts. The centre offset and the change in overlap shows strong sensitivity to
the centre definition. A single output and volume do not establish transfer
across all hydrodynamical models or redshifts.

The tabulated host and projection metrics, method summaries, paired
contrasts, metric definitions, position list, catalogue access and selection details, processing notes,
and source scripts appear in
\url{supplement/tng300_subhalopos_r50_transfer_20261103/}. The
original \texttt{tng300\_*.csv} files in the supplement root retain the
SubhaloCM-centred sensitivity results. Full reruns also require the original
TNG snapshot cutouts and project environment.

\section{TNG300-1 merger-tree growth analysis}
\label{sec:tng-growth-methods}

We retrieve the SubLink main-progenitor branch for each of the 40 TNG300-1
descendant central subhaloes described above \citep{rodriguezgomez2015}.
The target outputs are snapshots 84, 67, and 50, at redshifts 0.197284,
0.503048, and 0.997294. We retain the 36 branches whose subhalo remains
the FoF central at all three outputs. This sample cut uses catalogue
central status and does not use image measurements. The independent unit
is one descendant branch.

At each output, we compare image-based selections with the stars assigned
to the central subhalo. We use exact stellar-particle IDs from the TNG
catalogue to form this reference membership. Summing current particle
masses for these IDs reproduces the catalogue stellar mass
\texttt{SubhaloMassType[4]} to within $6.3\times10^{-5}$ dex for every
branch and output. We remove resolved satellite members using the same
mask as in the single-output TNG transfer test. Competing selections are
the isophotal region, friends-of-friends stellar groups with
$b_{\mathrm{link}}=0.10$, 0.12, or 0.20, a projected circle of radius
30 pkpc, and three-dimensional spheres using the Darwin-calibrated radii
of 80 pkpc in $V$ and 60 pkpc in $g$. The $b_{\mathrm{link}}=0.12$
($V$) and 0.10 ($g$) values are the low-redshift Darwin calibrations;
the other linking parameters are comparison settings. The sphere radii
were selected from a grid of 10, 15, 20, 30, 45, 60, 80, and 120 pkpc
by maximizing the median host-level, three-view light-weighted overlap
across the 12 Darwin groups at $z=0.190$. These full-sample radii are
held fixed in TNG and are not refit using TNG images.
We use current stellar-particle mass for selected stellar-mass sums.
Projected $R_{50}$ is the radius containing half the selected stellar
flux, with each particle's full flux assigned at its projected position.
As a structural-size sensitivity, we also calculate projected half-mass
radii by replacing luminosity weights with current stellar-particle mass.

For each physical output, the primary image uses that output's redshift
as the observer redshift. We also render each output at fixed
$z_{\rm obs}=0.197284$. These control images retain the output's stellar
positions, ages, and metallicities, while applying the same observer
distance, spectral redshift, surface-brightness dimming, angular scale,
and PSF. Comparing the two conditions measures how the adopted observer
redshift changes the result in this mock-image model; it does not
represent a separate physical history. The overlap measurements for all
three outputs and observer redshifts are tabulated. The main growth
comparison uses images at the physical redshifts and at fixed
$z_{\rm obs}=0.197284$.

For each branch and method, we calculate the growth bias in each of the
three orthogonal projections, then take the median across projections.
For every image, we record whether the centre pixel lies inside the
isophotal region, whether selection used the region containing the
nearest pixel brighter than the threshold, the region area, and whether
the region touches the image edge. This nearest-pixel selection is
retained and flagged separately from a central-pixel detection. The stellar-mass growth bias is
\begin{equation}
\Delta_{M}=\log_{10}\!\left[\frac{M_{\rm sel}(z=0.197)}{M_{\rm sel}(z=0.997)}\right]
-\log_{10}\!\left[\frac{M_{\rm cen}(z=0.197)}{M_{\rm cen}(z=0.997)}\right],
\end{equation}
where $M_{\rm sel}$ is the mass selected by a method and $M_{\rm cen}$ is
the mass in the central-subhalo reference. We define the projected-size
bias in the same way using $R_{50}$. Positive values mean that the method
measures more growth than the central-subhalo reference. We report the
sample median and percentile intervals from 20,000 bootstrap resamples of
the 36 descendant branches. We resample branches, not outputs or
projections.

The merger-tree tables, particle-membership checks, measurements for
each image, growth summaries, paired observer-redshift comparisons,
plotting and analysis scripts, and figure source are in
\url{supplement/tng300_growth_bias_20261106/}. Reproduction requires the
TNG300-1 cutouts and the stellar-population synthesis inputs documented in
that directory.

\begin{table*}[t]
\centering
\footnotesize
\setlength{\tabcolsep}{3pt}
\begin{tabular*}{\textwidth}{@{\extracolsep{\fill}}llrrrrrr@{}}
\toprule
Band & Selection & $\Delta_M$ & MAE$_M$ & $\Delta_{R_{50,\rm light}}$ & MAE$_{R_{50,\rm light}}$ & $\Delta_{R_{50,\rm mass}}$ & MAE$_{R_{50,\rm mass}}$ \\
\midrule
$V$ & Isophotal region & +0.176 & 0.184 & +0.430 & 0.519 & +0.315 & 0.316 \\
$V$ & FoF, $b_{\mathrm{link}}=0.10$ & -0.004 & 0.044 & +0.065 & 0.257 & -0.007 & 0.131 \\
$V$ & FoF, $b_{\mathrm{link}}=0.12$ & -0.003 & 0.037 & +0.011 & 0.235 & -0.004 & 0.112 \\
$V$ & FoF, $b_{\mathrm{link}}=0.20$ & -0.003 & 0.025 & -0.026 & 0.178 & +0.000 & 0.073 \\
$V$ & Projected aperture, 30 pkpc & -0.127 & 0.125 & -0.104 & 0.254 & -0.171 & 0.226 \\
$V$ & Spherical aperture, 80 pkpc & -0.067 & 0.081 & -0.147 & 0.218 & -0.125 & 0.150 \\
$g$ & Isophotal region & +0.261 & 0.531 & +0.296 & 0.468 & +0.323 & 0.424 \\
$g$ & FoF, $b_{\mathrm{link}}=0.10$ & -0.004 & 0.044 & +0.035 & 0.271 & -0.007 & 0.131 \\
$g$ & FoF, $b_{\mathrm{link}}=0.12$ & -0.003 & 0.037 & +0.025 & 0.258 & -0.004 & 0.112 \\
$g$ & FoF, $b_{\mathrm{link}}=0.20$ & -0.003 & 0.025 & -0.044 & 0.216 & +0.000 & 0.073 \\
$g$ & Projected aperture, 30 pkpc & -0.127 & 0.125 & -0.150 & 0.275 & -0.171 & 0.226 \\
$g$ & Spherical aperture, 60 pkpc & -0.088 & 0.098 & -0.166 & 0.267 & -0.142 & 0.180 \\
\bottomrule
\end{tabular*}
\caption{Bias in the growth of selected stellar mass and projected radius
relative to the central-subhalo reference, for images at the physical
observer redshift of each output. $\Delta_M$ is the median branch difference
in stellar-mass growth; the two $\Delta_{R_{50}}$ columns give the median
branch differences in half-light- and half-mass-radius growth. Each branch
contributes the median of three projections. The MAE columns give the mean
absolute branch difference, in dex, over 36 branches.}
\label{tab:growth-methods}
\end{table*}

\subsection{Detection, mass recovery, and size-weighting checks}
At the physical observer redshift of snapshot 50 ($z=0.997$), all 108
$V$ projections have a central-pixel detection. In $g$, 93/108 projections
have a central-pixel detection and 15 use the region containing the
nearest pixel above the threshold;
none has an empty region or a region touching the image edge. At snapshot 84
($z=0.197$), both bands have detections in all 108 projections. Requiring
central-pixel detections in all three views at both endpoints retains 36
branches in $V$ and 29 in $g$. The physical-redshift $g$-band mass-growth
bias is $0.261$ dex (95 per cent interval $0.182$--$0.303$) in the full
36-branch sample and $0.196$ dex ($0.148$--$0.289$) in the 29-branch
detection-qualified sample. The corresponding mean biases are $0.526$ and
$0.210$ dex. The full-sample estimate retains these nearest-pixel selections; the
separate detection-qualified estimate shows their effect without removing
them from the analysis.

The median selected-to-central mass ratios for $b_{\mathrm{link}}=0.10$
are $0.708$, $0.704$, and $0.708$ at snapshots 50, 67, and 84; for
$b_{\mathrm{link}}=0.20$ they are $0.852$, $0.836$, and $0.833$. The
corresponding tree-level 16th--84th percentile ranges are $0.62$--$0.78$
and $0.77$--$0.88$. For the calibrated $b_{\mathrm{link}}=0.12$ in $V$,
the median mass-growth bias is $-0.003$ dex and the mean absolute error is
$0.037$ dex. The calibrated $b_{\mathrm{link}}=0.10$ in $g$ gives
$-0.004$ dex and $0.044$ dex. The $b_{\mathrm{link}}=0.20$ comparison has
a mean absolute mass-growth error of $0.025$ dex in both bands.

The central-subhalo projected half-mass radius grows by a median of
$0.382$ dex (95 per cent interval $0.306$--$0.450$). The central-subhalo
half-light-radius growth is $0.020$ dex in $V$ and $-0.003$ dex in $g$
for physical-redshift images; both intervals include zero. At fixed observer
redshift, the half-light-radius growth medians are $0.257$ dex in $V$ and
$0.237$ dex in $g$. Thus the half-light-radius result depends on observer
condition and luminosity weighting, while the half-mass radius provides a
separate structural comparison.

The R1--R4 results directory contains central-pixel detections and
nearest-pixel selections for each branch and projection, isophotal areas
and edge flags, stellar-mass recovery ratios, both growth samples,
half-mass-radius measurements, bootstrap intervals, and the updated
figure source. The former
growth summaries in the parent directory use only $b_{\mathrm{link}}=0.10$
and 0.20 and are superseded for the revised comparison.

The branch-level mass-growth errors show a broad positive tail for the
isophotal selection. The 16th--84th percentile range is $0.077$--$0.278$
dex in $V$ and $0.069$--$0.691$ dex in $g$ for all 36 branches. Among the
29 $g$-band branches with central-pixel detections in every view at both
endpoints, the range is $0.054$--$0.323$ dex. By comparison, the calibrated
FoF growth errors remain close to zero. The seven branches excluded by the
$g$-band detection cut include the largest positive errors; the full-sample
median retains these cases rather than discarding them.

The nearly constant FoF mass-recovery ratios explain why substantial mass
under-recovery can coexist with a small growth error: a ratio that changes
little between outputs largely cancels in the growth comparison. The
isophotal recovery rises from $0.479$ to $0.747$ in $V$ and from $0.406$ to
$0.719$ in $g$ between snapshots 50 and 84, so its growth error is more
positive. Figure~\ref{supp:branchgrowth} shows the branch-level errors and
the recovery distributions at all three outputs.

\begin{figure*}
\centering
\includegraphics[width=\textwidth]{supplement/astra_review_response_20261106/branch_level_growth_recovery.pdf}
\caption{Branch-level stellar-mass growth error and mass recovery in
TNG300-1. Panels (a) and (b) show the change in selected stellar-mass growth
minus the central-subhalo growth from snapshot 50 to 84. Each point is one
branch after taking its median over three projections; horizontal marks show
the sample median. Filled circles mark branches with central-pixel
detections in all three views at both endpoints. Open triangles mark the
seven $g$-band branches with at least one nearest-pixel fallback. The inset
in panel (b) expands the values near zero. Panels (c) and (d) show the
selected-to-central stellar-mass ratio at snapshots 50, 67, and 84. Curves
are medians across branches and shaded regions span the 16th--84th
percentiles. The linking parameters in the calibrated FoF curves are 0.12
in $V$ and 0.10 in $g$; the comparison curve uses 0.20.}
\label{supp:branchgrowth}
\end{figure*}

\subsection{Stellar-mass budget and stellar-population controls}
\label{sec:tng-assembly-population}

We use the exact \texttt{ParticleIDs} in the central-subhalo catalogues at
snapshots 50, 67, and 84 to separate changes in central membership from
changes in particle mass. For each pair of outputs, a particle is persistent
if it belongs to the central subhalo at both outputs. Pre-existing newly
central stars have formation scale factors at or before the first output and
are central only at the second output. Newly formed central stars formed
between the outputs and are central at the second output. The lost category
contains particles central at the first output that are no longer central at
the second. We retain particles with finite stellar-formation scale factors
$0<a_{\rm birth}\leq a_{\rm snap}$; the excluded counts are listed in
\texttt{input\_membership\_integrity.csv}. We use current particle masses in
the mass balance. The equation is
\begin{equation}
\Delta M_{\rm cen}=\Delta M_{\rm persistent}
+M_{\rm old,new\ cen}
+M_{\rm formed,new\ cen}
-M_{\rm lost\ cen}.
\end{equation}
We evaluate the mass balance for each branch before summarizing the sample. We
normalize each term by the central stellar mass at the starting output.
The maximum absolute closure residual over all 36 branches and all three
output pairs is $9.7\times10^{-16}$ of the starting mass.

To characterize the earlier catalogue assignment of pre-existing newly
central stars, we match their IDs to the earlier parent-FoF cutout and record
the \texttt{RemoveSatelliteGalaxy} flag. That flag marks stars assigned to
resolved non-central subhaloes with at least 100 stellar particles; the
cutouts do not provide the source subhalo ID for these stars. Stars absent
from the earlier parent-FoF cutout have an unknown earlier location. We
therefore report catalogue-mask membership and do not interpret it as a
measured merger or accretion history. We also do not label stars formed
between outputs as in situ because their birth hosts are not recorded.

The size controls use the same central-subhalo members, proper positions,
and dust-free FSPS simple-stellar-population flux grids as the reference
half-light radii. In the controls, we replace each particle's age, total
metal mass fraction, or both with the central population's geometric mean
for that quantity at that output. Each geometric mean is the exponential
of the current-mass-weighted mean logarithm. This changes the spatial population
weighting while leaving positions and formed-mass normalizations unchanged.
The age and metallicity contributions use a two-factor Shapley split, which
assigns half of their interaction to each factor. The combined effect is
computed as a paired full-population minus uniform-population growth for
each branch and projection before reducing projections to a branch median.
The statistical unit is the descendant branch; 95 per cent intervals use
20,000 bootstrap resamples of the 36 branches. Radial profiles combine bins
inside and outside one projected current-mass half-mass radius. The quoted
metallicity is the total metal mass fraction in \texttt{GFM\_Metallicity},
not [Fe/H]; \texttt{GFM\_StellarFormationTime} records the formation scale
factor [\url{https://www.illustris-project.org/data/docs/specifications/}].

From snapshot 50 to 84, the median net central stellar-mass increase is
$0.869$ times the starting central mass. The median mass of pre-existing
stars newly assigned to the central is $0.789$ times the starting mass;
the median current-mass change of persistent central members is $-0.050$,
and stars formed after snapshot 50 and central at snapshot 84 contribute
$0.096$. The mass of snapshot-50 central particles no longer central at
snapshot 84 contributes $-0.002$. These are separate component medians and
do not sum to the median net change. A median 88.1 per cent of the pre-existing
newly central mass carried the earlier resolved-satellite flag. The unflagged
remainder is 11.9 per cent in the branch-level medians. Of the total newly
central pre-existing mass, 2.1 per cent was in the earlier parent FoF without
the satellite flag, while 8.2 per cent was absent from that earlier FoF
catalogue. The component medians do not add to the median remainder; the
earlier location of absent stars is unknown.
Table~\ref{tab:tng-id-budget} gives the bootstrap intervals across branches.

\begin{table*}
\centering
\caption{Particle-ID mass budget for the change in central-subhalo stellar mass
from snapshot 50 ($z=0.997$) to snapshot 84 ($z=0.197$). The first five
quantities are normalized by each branch's snapshot-50 central stellar mass;
the final quantity is the fraction of pre-existing newly central mass with
the earlier resolved-satellite flag. Values give the sample median and the
95 per cent branch-bootstrap interval. Component medians do not add because
the mass balance closes separately for each branch.}
\label{tab:tng-id-budget}
\small
\begin{tabular}{lr}
\toprule
Mass-budget component & Median [95 per cent interval] \\
\midrule
Net central stellar-mass increase / $M_{\rm cen,50}$ & $+0.869\;[+0.650,+1.332]$ \\
Persistent-particle current-mass change / $M_{\rm cen,50}$ & $-0.050\;[-0.054,-0.046]$ \\
Pre-existing stars newly assigned to central / $M_{\rm cen,50}$ & $+0.789\;[+0.631,+0.971]$ \\
Stars formed after snapshot 50, central at 84 / $M_{\rm cen,50}$ & $+0.096\;[+0.060,+0.348]$ \\
Snapshot-50 central stars no longer central at 84 / $M_{\rm cen,50}$ & $-0.002\;[-0.003,-0.001]$ \\
Earlier resolved-satellite-flag fraction of pre-existing newly central mass & $0.881\;[0.766,0.919]$ \\
\bottomrule
\end{tabular}
\end{table*}

At snapshot 84, the linked groups and isophotal regions capture different
fractions of these particle populations (Table~\ref{tab:tng-capture}). The
FoF values are independent of photometric band because the 3D membership is
the same in $V$ and $g$. Increasing $b_{\mathrm{link}}$ from 0.10 to 0.20
raises capture of resolved-satellite stars and newly formed stars, while
also including more stars that were not central members and did not carry the
earlier satellite flag.

\begin{table*}
\centering
\caption{Fraction of each snapshot-84 stellar population selected by the
isophotal region or linked group. Values are percentages: the three
projections are reduced to a median for each branch, followed by the median
across 36 branches. The satellite class uses the earlier
\texttt{RemoveSatelliteGalaxy} flag. The other parent-FoF class was in the
earlier group but was neither central nor flagged as a resolved satellite;
the absent class has no matching star in that earlier FoF catalogue. The FoF
rows are identical in $V$ and $g$ and are shown once.}
\label{tab:tng-capture}
\small
\begin{tabular}{lrrrrr}
\toprule
Selection & Persistent & Earlier satellite & Other in parent FoF & Absent from parent FoF & Formed since 50 \\
\midrule
Isophotal, $V$ & 81.4 & 73.3 & 12.7 & 2.7 & 56.9 \\
Isophotal, $g$ & 79.1 & 70.3 & 11.0 & 2.3 & 52.8 \\
FoF, $b_{\mathrm{link}}=0.10$ & 78.1 & 69.9 & 7.1 & 0.8 & 52.8 \\
FoF, $b_{\mathrm{link}}=0.12$ & 80.9 & 74.3 & 9.5 & 1.5 & 59.5 \\
FoF, $b_{\mathrm{link}}=0.20$ & 88.8 & 84.1 & 19.0 & 4.7 & 71.6 \\
\bottomrule
\end{tabular}
\end{table*}

With the physical observer redshift at each output, the full-population
half-light-radius growth is $0.020$ dex in $V$ and $-0.003$ dex in $g$;
the uniform-age-and-metallicity control grows by $0.386$ dex in both bands.
This radius weights particles by formed stellar mass; the half-mass radius
uses current particle mass and grows by $0.382$ dex.
The paired full-minus-uniform change is $-0.255$ dex in $V$ and $-0.324$
dex in $g$. At fixed $z_{\rm obs}=0.197284$, the corresponding paired
changes are $-0.071$ dex in $V$ and $-0.085$ dex in $g$. The paired
full-population half-light-radius growth is lower by $0.183$ dex in $V$
and $0.199$ dex in $g$ at the physical observer redshifts than at fixed
$z_{\rm obs}$. For the full-minus-uniform growth change, the age and
metallicity Shapley contributions at the physical observer redshifts are
$-0.162$ and $-0.123$ dex in $V$, and $-0.175$ and $-0.143$ dex in $g$.
At fixed observer redshift, they are $-0.047$ and $-0.018$ dex in $V$, and
$-0.051$ and $-0.007$ dex in $g$. The contribution medians need not sum to
the median combined effect; branch-bootstrap intervals are in
\texttt{results/key\_results.csv}.

At snapshot 50, stars outside one projected central half-mass radius are
median $0.304$ Gyr younger and lower by $0.514$ dex in log total metal mass
fraction than stars inside it. Their outer-minus-inner log mass-to-light
ratio is $-0.556$ dex in $V$ and $-0.775$ dex in $g$ when each output is rendered at
its physical observer redshift. At fixed $z_{\rm obs}=0.197284$, the
snapshot-50 differences are $-0.245$ dex in $V$ and $-0.301$ dex in $g$.
At snapshot 84, the outer-minus-inner differences are $-0.558$ Gyr,
$-0.435$ dex in log metallicity, and $-0.168$ dex in $V$ and $-0.205$ dex
in $g$ under either observer condition. The stronger snapshot-50 contrast
at its physical observer redshift occurs when the observed bands sample
rest-frame near-UV light. These radial contrasts describe the simulated
population gradients. The controlled SSP comparison uses dust-free spectra;
dust attenuation could change the measured contrast, and we do not estimate
that change here. Per-projection gradients and
all branch-bootstrap summaries are in
\texttt{results/radial\_gradients\_per\_view.csv} and
\texttt{results/radial\_gradient\_summary.csv}.

\section{Artificial-redshift sample and detection accounting}
The experiment holds each halo fixed in proper size at the original simulation
output and renders all 75 entries from outputs $z\leq1.990$ at observed
redshifts 0.190, 0.227, 0.500, 1.002, and 1.990. We convert the stored
comoving coordinates to proper positions using the simulation scale factor,
then use the angular-diameter distance at the observed redshift. The angular
pixel scale and PSF stay fixed. The physical resolution therefore varies
with redshift. Changes combine
cosmological dimming, spectral shifting, and physical resolution, and
do not isolate the causal effect of any one ingredient. The four flux
modes are defined as follows. \texttt{observer\_full} uses the redshifted
stellar spectrum and corrects the FSPS flux normalization to the Darwin
luminosity distance. \texttt{observer\_no\_tolman\_factor} multiplies
the corrected flux by $(1+z_{\rm obs})^4$. \texttt{rest\_sed\_with\_distance}
uses the $z=0$ stellar spectrum with the Darwin luminosity distance.
\texttt{rest\_sed\_no\_dimming} uses the $z=0$ absolute spectrum
without luminosity-distance scaling. The final mode is an extreme no-distance control,
not merely the removal of the Tolman factor.

The data table \texttt{artificial\_redshift\_summary.csv} gives the fraction
of views with a detected isophote, the mean overlap among detected views,
and the mean overlap across all views. Every view has equal weight. Views
without a detected isophote have zero overlap in the all-view mean.
Completeness is undefined without a detected region and is averaged only over
detections. The table reports the difference between the all-view mean and
the detection fraction multiplied by the mean overlap among detections.
The difference is zero to numerical precision.

At $z_{\rm obs}=1.990$, isophotes are detected in 182/225 $V$ views and
188/225 $g$ views. The view-weighted mean overlap among detected views is
0.829 in $V$ and 0.860 in $g$. Mean completeness among detected views is
0.996 and 0.994. Assigning zero overlap to non-detections gives means of
0.671 and 0.719 across all views. The means among detections describe
images where an isophote was found, while the all-view means also account
for views without a detected isophote. Removing the $(1+z)^4$ cosmological-dimming factor raises the
all-view means by 0.211 in $V$ and 0.158 in $g$. Using the $z=0$ stellar
spectrum while retaining the Darwin luminosity distance gives means of
0.795 and 0.818. The extreme control without luminosity-distance scaling is limited by
the finite 360-arcsec field. The full stamp is filled in 81, 67, 15, 3,
and 2 of 225 views at observed redshifts 0.190, 0.227, 0.500, 1.002,
and 1.990, respectively, in each band. The mean overlaps are therefore
field-of-view dependent and are not interpreted as physical comparisons.
The controls change image construction in different ways and do not
isolate one physical cause or represent evolution of the stellar systems.

An additional comparison pairs image conditions within the 33 major exact-ID
components. The calculation first takes the median over three projections at
each halo output and then takes the median over repeated outputs in each
component. Removing the $(1+z_{\rm obs})^4$ surface-brightness dimming lowers
mean component IoU at $z_{\rm obs}=0.190$ by $0.0273$ in $V$ and $0.0235$ in
$g$. At $z_{\rm obs}=1.990$, removing dimming raises mean component IoU by
$0.0834$ in $V$ and $0.0485$ in $g$. The 95 per cent interval includes zero
for the $g$-band change. Replacing redshifted stellar spectra with $z=0$
spectra raises mean IoU at $z_{\rm obs}=1.990$ by $0.0958$ in $V$ and
$0.0769$ in $g$. The paired changes vary with observed redshift and band.
The angular pixel and PSF scales remain fixed, so physical resolution changes
with observed redshift. The paired changes measure image-construction
sensitivity and do not measure halo evolution.

\begin{table*}
\centering
\caption{Paired changes in overlap for artificial-redshift image controls at
$b_{\mathrm{link}}=0.10$. Each value gives the mean component-level change
from the full image with redshifted stellar spectra, followed by the 95 per
cent bootstrap interval. The first control removes cosmological surface-brightness
dimming while retaining the redshifted spectra and luminosity distance. The
second control uses the $z=0$ stellar spectra and retains the luminosity
distance. The comparison uses 53 halo-output entries from outputs 00178,
00475, and 00786. Twenty major-descendant exact dark-matter particle-ID links
form 33 connected components. The analysis takes the median across three
projections at each output and then the median across repeated outputs in each
component. Bootstrap resampling uses the 33 components. The components do not
represent a complete merger-tree census, and image-mode contrasts also include
the change in physical resolution with observed redshift.}
\label{tab:redshiftmodepaired}
\small
\setlength{\tabcolsep}{4pt}
\begin{tabular}{ccccc}
\toprule
$z_{\rm obs}$ & Band & $N_{\rm comp}$ & No-dimming $\Delta\mathrm{IoU}$
& $z=0$ spectrum $\Delta\mathrm{IoU}$ \\
\midrule
0.190 & $V$ & 33 & $-0.0273\;[-0.0334,-0.0216]$ & $-0.0070\;[-0.0087,-0.0052]$ \\
 & $g$ & 33 & $-0.0235\;[-0.0299,-0.0179]$ & $-0.0161\;[-0.0195,-0.0129]$ \\
0.227 & $V$ & 33 & $-0.0305\;[-0.0377,-0.0242]$ & $-0.0081\;[-0.0101,-0.0062]$ \\
 & $g$ & 33 & $-0.0250\;[-0.0322,-0.0184]$ & $-0.0165\;[-0.0197,-0.0133]$ \\
0.500 & $V$ & 33 & $-0.0176\;[-0.0294,-0.0063]$ & $+0.0003\;[-0.0061,+0.0070]$ \\
 & $g$ & 33 & $-0.0086\;[-0.0202,+0.0027]$ & $-0.0011\;[-0.0094,+0.0084]$ \\
1.002 & $V$ & 33 & $+0.0264\;[-0.0010,+0.0604]$ & $+0.0389\;[+0.0068,+0.0778]$ \\
 & $g$ & 33 & $+0.0201\;[-0.0124,+0.0643]$ & $+0.0379\;[-0.0031,+0.0900]$ \\
1.990 & $V$ & 33 & $+0.0834\;[+0.0167,+0.1688]$ & $+0.0958\;[+0.0162,+0.1901]$ \\
 & $g$ & 33 & $+0.0485\;[-0.0076,+0.1192]$ & $+0.0769\;[+0.0030,+0.1657]$ \\
\bottomrule
\end{tabular}
\end{table*}

\begin{table*}
\centering
\caption{Fixed-$b_{\mathrm{link}}$ comparison after grouping repeated outputs
with major exact-ID links. The sample contains 53 entries from outputs 00178,
00475, and 00786. Twenty major-descendant links define 33 components. Each
component first takes the median across three views at each output, followed
by the median across repeated outputs. $\Delta\mathrm{IoU}$ gives the mean
paired difference $b_{\mathrm{link}}=0.10$ minus $b_{\mathrm{link}}=0.20$ and
the 95 per cent interval from 5,000 bootstrap resamples of the 33 components.
The failure fraction counts components with median IoU below 0.90 at the two
fixed linking lengths. $P_{90}(|\Delta\log_{10}L|)$ gives the 90th-percentile
absolute luminosity error in dex at $b_{\mathrm{link}}=0.10$ and 0.20. IoU
includes zero for views without a detected isophote. Luminosity errors use
components with a detected isophote. At $z_{\rm obs}=1.990$, 30 $V$ components
and 32 $g$ components have valid luminosity errors for each linking length.
All other luminosity-error quantiles use 33 components.}
\label{tab:exactredshift}
\small
\setlength{\tabcolsep}{5pt}
\begin{tabular}{ccccc}
\toprule
$z_{\rm obs}$ & Band & $\Delta\mathrm{IoU}$ (95\% interval)
& $f_{<0.90}$ (\%, $b_{\mathrm{link}}=0.10/0.20$)
& $P_{90}(|\Delta\log_{10}L|)$ (dex, $b_{\mathrm{link}}=0.10/0.20$) \\
\midrule
0.190 & $V$ & $-0.0105\;[-0.0224,+0.0001]$ & 15/9 & 0.049/0.027 \\
 & $g$ & $-0.0064\;[-0.0180,+0.0045]$ & 15/9 & 0.049/0.032 \\
0.227 & $V$ & $-0.0037\;[-0.0144,+0.0065]$ & 12/9 & 0.044/0.032 \\
 & $g$ & $+0.0021\;[-0.0081,+0.0124]$ & 12/12 & 0.043/0.039 \\
0.500 & $V$ & $+0.0388\;[+0.0283,+0.0498]$ & 15/33 & 0.024/0.076 \\
 & $g$ & $+0.0313\;[+0.0223,+0.0407]$ & 15/30 & 0.029/0.089 \\
1.002 & $V$ & $+0.0395\;[+0.0266,+0.0539]$ & 21/27 & 0.120/0.238 \\
 & $g$ & $+0.0286\;[+0.0144,+0.0445]$ & 18/27 & 0.134/0.185 \\
1.990 & $V$ & $+0.0206\;[+0.0110,+0.0329]$ & 24/27 & 0.067/0.123 \\
 & $g$ & $+0.0221\;[+0.0107,+0.0355]$ & 21/27 & 0.144/0.253 \\
\bottomrule
\end{tabular}
\end{table*}

\section{Fixed-parameter image checks in CAMELS}
\label{sec:camelsimage}
The image checks use the same 49 hosts, potential-minimum centres,
stellar masks, observing redshift, and corrected flux convention as the
nominal CAMELS transfer test. The linking parameters remain fixed at
$b_{\mathrm{link}}=0.10$ and $0.20$. A projected 30 pkpc aperture provides a third
selection. The checks were designed after the nominal transfer result
and do not select a new linking parameter from the external sample.

The baseline uses $0.2$-arcsec pixels, a $1$-arcsec FWHM Moffat PSF
with index $4.765$ and the outer wing defined in
Section~\ref{sec:psfwing}, and limits
of $26.5$ and $27.0\,\mathrm{mag\,arcsec}^{-2}$ in $V$ and $g$.
Each check changes one choice. The limits become $0.5$ mag brighter
or fainter, the pixels become $0.4$ arcsec wide, the PSF wing is
removed, or a constant positive residual of
$29\,\mathrm{mag\,arcsec}^{-2}$ is added to the image. The no-wing
PSF is normalized to unit flux. The field of view remains fixed.
Three further checks add independent, zero-mean Gaussian noise after
PSF convolution. The pixel rms equals the pixel flux at the brightness
limit divided by 5, 3, or 1. Negative pixel values remain in the
noisy images. Each host, view, and band has three deterministic noise
realisations. The three noise amplitudes use the same standard-normal
patterns, scaled to the requested rms. The Gaussian model does not
represent a calibrated survey noise distribution.

For each realisation, we take the median over three views. We then
average the three realisations within each host. Noiseless conditions
use one realisation. Equal-host means and 90th-percentile absolute
errors describe the 49 host summaries. The paired differences use
5,000 shared bootstrap draws, resampling hosts or resampling boxes
and then hosts within each sampled box. Views and noise realisations
do not increase the number of independent hosts. An absent central region or a region selected around the nearest pixel
above the threshold counts as a failed target. Failed targets receive zero overlap and
completeness. Precision and luminosity or radius errors are undefined
for failed targets. Error quantiles therefore condition on measurable
targets and failure fractions are reported separately.

Every host reproduces the nominal baseline projection metrics before
entering the summary. The stronger-noise extension also reproduces the
baseline. The combined projection table removes the repeated baseline
and contains 13,230 rows across nine image conditions, two bands, and
three stellar selections. No central target fails in any condition.
The mean overlap difference between $b_{\mathrm{link}}=0.10$ and $0.20$ has a positive
95 per cent box-then-host interval in every condition and band. At the
largest noise amplitude, the differences are $+0.075$ in $V$
($+0.063$--$+0.088$) and $+0.090$ in $g$
($+0.075$--$+0.107$). The mean overlaps at $b_{\mathrm{link}}=0.10$ fall from
$0.894/0.872$ to $0.877/0.857$ ($V/g$).

Changing the brightness limit shifts the luminosity bias more than
changing the pixel size or removing the PSF wing. With fainter limits,
the $V$ mean overlap rises to $0.901$, while the 90th-percentile
absolute luminosity error rises from $0.033$ to $0.056$ dex.
The isophotal target changes with every image perturbation. A higher
overlap does not guarantee a smaller luminosity error. Full noiseless
particle fluxes at projected positions define luminosity and $R_{50}$.
The image checks measure sensitivity to the boundary selection rather
than noise in the measured stellar flux. Dust, correlated sky noise,
Poisson counting statistics, and image-based centre finding are omitted.
Figure~\ref{fig:camelsimage} compares overlap and luminosity errors.
The directory \texttt{camels\_image\_checks} supplies the projection
table, host summaries, paired differences, configurations, and
provenance. The script
\texttt{summarize\_camels\_image\_checks.py} reproduces the summary
tables without the original particle snapshots.

\begin{figure*}
\centering
\includegraphics[width=\textwidth]{supplement/camels_image_checks/fixed_b_image_sensitivity.pdf}
\caption{Fixed-parameter image checks for 49 CAMELS Swift-EAGLE hosts.
The upper row shows equal-host mean overlaps, and the lower row shows
the 90th percentile of the absolute logarithmic luminosity error in dex,
$|\Delta\log_{10}L|$, where
$\Delta\log_{10}L=\log_{10}(L_{\rm selected}/L_{\rm isophote})$.
The columns show $V$ and SDSS $g$. Base denotes the nominal image.
The next conditions change the brightness limit by $-0.5$ mag
(brighter) or $+0.5$ mag (fainter), use $0.4$-arcsec pixels, remove
the outer PSF wing, or add a flat positive residual of
$29\,\mathrm{mag\,arcsec}^{-2}$. The final conditions add Gaussian
pixel noise with rms equal to the threshold pixel flux divided by
5, 3, or 1. Each noisy host statistic averages the three-view medians
from three noise realisations. The linking parameters remain fixed.
The 30 pkpc aperture uses the same known centre. All luminosities
use noiseless stellar-particle fluxes.}
\label{fig:camelsimage}
\end{figure*}

\section{Fixed-structure redshift transfer in CAMELS}
\label{sec:camelsmultiz}

We render the 49 CAMELS Swift-EAGLE hosts from snapshot 090 at four observer
redshifts: $z_{\rm obs}=0.190$, $0.500$, $1.002$, and $1.990$. Stellar
positions, ages, metallicities, and FoF membership remain fixed at the $z=0$
snapshot. The image calculation changes the redshifted stellar spectra,
luminosity and angular-diameter distances, cosmological surface-brightness
dimming, and physical sampling of the fixed angular pixel and PSF. Each
redshift uses the nominal limits of $26.5$ and
$27.0\,\mathrm{mag\,arcsec}^{-2}$ in $V$ and $g$, plus four sensitivity
cases that make one band limit $0.5$ mag brighter or fainter. We keep
$b_{\mathrm{link}}=0.10$, $b_{\mathrm{link}}=0.20$, and the projected 30 pkpc aperture
fixed in every case. The calculation omits dust, sky background, pixel noise,
and intergalactic absorption. Some $z=0$ stellar ages exceed the age of the
Universe at $z_{\rm obs}=1.990$. These images therefore test how the
measurement changes with observer redshift; they are not evolving galaxy
models.

For each host and selection, we first take the median over its three
orthogonal projections. The mean overlap then gives equal weight to all 49
hosts. A missing isophotal target in any view contributes zero overlap to
that host's view median. $N_{\rm fail}$ counts hosts with at least one failed
view. The paired intervals resample the seven CV boxes and then hosts within
each selected box. Views are repeated measurements, not independent hosts.
Table~\ref{tab:camelsmultizbase} gives the nominal-limit results for
$b_{\mathrm{link}}=0.10$. Figure~\ref{fig:camelsmultiz} compares the three fixed
selections across observer redshift.

At the nominal limits, $b_{\mathrm{link}}=0.10$ has a higher mean overlap than
$b_{\mathrm{link}}=0.20$ at all four redshifts, but the absolute overlap falls
sharply above $z_{\rm obs}=0.190$. At $z_{\rm obs}=0.500$ and $1.002$, the
30 pkpc aperture has a higher mean overlap than the fixed linking parameter.
At $z_{\rm obs}=1.002$ and $1.990$, many hosts have at least one failed view.
The luminosity, size, and contamination comparisons at those redshifts
therefore use only hosts with measurable isophotes and do not represent the
full sample. Changing a limit by $0.5$ mag changes the overlap and failure
count, as listed in Table~\ref{tab:camelsmultizthreshold}. In the deeper
single-band cases at $z_{\rm obs}=1.990$, the affected band still has failed
views in 43 $V$ hosts or 41 $g$ hosts. The threshold comparison does not
separate dimming from spectral shifting or physical resolution.

\begin{figure*}
\centering
\includegraphics[width=\textwidth]{supplement/camels_multiz_20261030/camels_multiz_transfer_summary_v2.pdf}
\caption{Fixed-structure image-transfer test for 49 CAMELS Swift-EAGLE
hosts. The top panels show mean host-median overlap at the nominal
surface-brightness limits. The bottom panels show the fraction of hosts
with at least one failed projection; the dashed grey line is the same for
all three selections because contour detection depends on the image. Error
bars in the top panels show 95 per cent intervals from resampling CV boxes
and hosts within boxes. Stellar structure and FoF membership are fixed
while observer redshift changes.}
\label{fig:camelsmultiz}
\end{figure*}

\begin{table*}
\centering
\caption{Nominal-limit redshift-transfer results for fixed
$b_{\mathrm{link}}=0.10$. Mean IoU is the equal-host mean of each host's
median over three projections. $N_{\rm IoU<0.90}$ counts hosts with a
host-median IoU below 0.90. $N_{\rm fail}$ counts hosts with at least one
projection lacking a measurable isophote. The two $\Delta$IoU columns give
paired mean differences and 95 per cent box-then-host bootstrap intervals.}
\label{tab:camelsmultizbase}
\footnotesize
\setlength{\tabcolsep}{3pt}
\begin{tabular}{ccrrrll}
\toprule
$z_{\rm obs}$ & Band & Mean IoU & $N_{\rm IoU<0.90}$ & $N_{\rm fail}$
& $\Delta$IoU vs $b_{\mathrm{link}}=0.20$
& $\Delta$IoU vs 30 pkpc \\
\midrule
0.190 & $V$ & 0.894 & 32 & 0 & $+0.094\;[+0.083,+0.106]$ & $+0.040\;[+0.005,+0.083]$ \\
      & $g$ & 0.872 & 40 & 0 & $+0.111\;[+0.098,+0.127]$ & $+0.010\;[-0.025,+0.051]$ \\
0.500 & $V$ & 0.690 & 47 & 0 & $+0.121\;[+0.105,+0.140]$ & $-0.130\;[-0.173,-0.099]$ \\
      & $g$ & 0.653 & 46 & 0 & $+0.128\;[+0.110,+0.149]$ & $-0.125\;[-0.172,-0.090]$ \\
1.002 & $V$ & 0.263 & 46 & 29 & $+0.047\;[+0.025,+0.070]$ & $-0.037\;[-0.059,-0.018]$ \\
      & $g$ & 0.300 & 43 & 32 & $+0.054\;[+0.028,+0.082]$ & $-0.041\;[-0.067,-0.019]$ \\
1.990 & $V$ & 0.041 & 49 & 46 & $+0.004\;[+0.000,+0.017]$ & $-0.004\;[-0.014,+0.000]$ \\
      & $g$ & 0.116 & 49 & 44 & $+0.022\;[+0.002,+0.048]$ & $-0.013\;[-0.030,+0.000]$ \\
\bottomrule
\end{tabular}
\end{table*}

\begin{table*}
\centering
\caption{Surface-brightness-limit sensitivity at fixed
$b_{\mathrm{link}}=0.10$. Each row uses all 49 hosts. The limit pair is
$(\mu_{V,\lim},\mu_{g,\lim})$ in $\mathrm{mag\,arcsec}^{-2}$. Mean IoU
includes all hosts, assigning zero overlap when a target is not measured.
$N_{\rm fail}$ counts hosts with at least one failed projection.}
\label{tab:camelsmultizthreshold}
\small
\setlength{\tabcolsep}{5pt}
\begin{tabular}{ccrrrr}
\toprule
$z_{\rm obs}$ & $(\mu_{V,\lim},\mu_{g,\lim})$ & Mean IoU ($V$) & $N_{\rm fail}(V)$
& Mean IoU ($g$) & $N_{\rm fail}(g)$ \\
\midrule
0.190 & $(26.5,27.0)$ & 0.894 & 0 & 0.872 & 0 \\
      & $(26.0,27.0)$ & 0.858 & 0 & 0.872 & 0 \\
      & $(27.0,27.0)$ & 0.901 & 0 & 0.872 & 0 \\
      & $(26.5,26.5)$ & 0.894 & 0 & 0.827 & 0 \\
      & $(26.5,27.5)$ & 0.894 & 0 & 0.892 & 0 \\
0.500 & $(26.5,27.0)$ & 0.690 & 0 & 0.653 & 0 \\
      & $(26.0,27.0)$ & 0.625 & 0 & 0.653 & 0 \\
      & $(27.0,27.0)$ & 0.748 & 0 & 0.653 & 0 \\
      & $(26.5,26.5)$ & 0.690 & 0 & 0.589 & 0 \\
      & $(26.5,27.5)$ & 0.690 & 0 & 0.713 & 0 \\
1.002 & $(26.5,27.0)$ & 0.263 & 29 & 0.300 & 32 \\
      & $(26.0,27.0)$ & 0.185 & 37 & 0.300 & 32 \\
      & $(27.0,27.0)$ & 0.379 & 10 & 0.300 & 32 \\
      & $(26.5,26.5)$ & 0.263 & 29 & 0.215 & 36 \\
      & $(26.5,27.5)$ & 0.263 & 29 & 0.392 & 24 \\
1.990 & $(26.5,27.0)$ & 0.041 & 46 & 0.116 & 44 \\
      & $(26.0,27.0)$ & 0.015 & 49 & 0.116 & 44 \\
      & $(27.0,27.0)$ & 0.134 & 43 & 0.116 & 44 \\
      & $(26.5,26.5)$ & 0.041 & 46 & 0.038 & 47 \\
      & $(26.5,27.5)$ & 0.041 & 46 & 0.165 & 41 \\
\bottomrule
\end{tabular}
\end{table*}

The tabulated summaries in
\texttt{camels\_multiz\_20261030/condition\_metrics.csv} give the 120
condition, method, and band summaries. Paired redshift and threshold
contrasts are in \texttt{paired\_redshift\_changes.csv} and
\texttt{paired\_threshold\_changes.csv}; the rendering configuration and
input checks are recorded in \texttt{summary\_metadata.json}.

\section{Archived image-construction and sampling checks}
The exploratory checks use the archived photometric normalization. The main
conclusions and quoted calibration use the corrected-distance calculations.
\begin{table*}
\centering
\caption{Sensitivity to image construction.
$N$ counts haloes retained for paired changes, and $P_{16}$--$P_{84}$
is the percentile interval of $\Delta b_{\mathrm{link}}$.
Edge counts views moving to either grid endpoint.
Below $0.90$ counts views remaining interior but falling below the
baseline overlap threshold.
The counts use the projection sample defined in the text, including qualifying
views of the excluded halo.
Sky brightnesses are in $\mathrm{mag\,arcsec}^{-2}$.
Ellipses denote the absence of a valid selected linking parameter.}
\label{tab:imagingvariants}
\small
\setlength{\tabcolsep}{5pt}
\begin{tabular}{llrrrrrr}
\toprule
Change & Band & $N$ & Median $\Delta b_{\mathrm{link}}$ & $P_{16}$--$P_{84}$
& Median $\Delta\mathrm{IoU}$ & Edge & Below $0.90$ \\
\midrule
Remove PSF wing & $V$ & 75 & $0$ & $[0,0]$ & $0$ & 0 & 0 \\
 & $g$ & 75 & $0$ & $[0,0]$ & $0$ & 0 & 0 \\
Sky $29$ & $V$ & 75 & $0$ & $[0,0.02]$ & $0.0008$ & 0 & 2 \\
 & $g$ & 74 & $0$ & $[0,0.02]$ & $0.0013$ & 1 & 0 \\
Sky $27$ & $V$ & 56 & $0.03$ & $[0.02,0.05]$ & $0.0066$ & 28 & 4 \\
 & $g$ & 0 & $\ldots$ & $\ldots$ & $\ldots$ & 198 & 0 \\
Threshold $0.5$ mag brighter & $V$ & 74 & $0$ & $[-0.02,0]$ & $-0.005$ & 1 & 15 \\
 & $g$ & 75 & $-0.02$ & $[-0.03,0]$ & $-0.0039$ & 0 & 11 \\
Threshold $0.5$ mag fainter & $V$ & 73 & $0.02$ & $[0,0.03]$ & $0.003$ & 2 & 3 \\
 & $g$ & 68 & $0.02$ & $[0,0.03]$ & $0.0024$ & 7 & 2 \\
Include satellite light & $V$ & 75 & $0$ & $[0,0.02]$ & $0$ & 1 & 5 \\
 & $g$ & 73 & $0$ & $[0,0.0096]$ & $0$ & 2 & 4 \\
\bottomrule
\end{tabular}
\end{table*}

\begin{table*}
\centering
\caption{Sensitivity to particle sampling.
The retained fractions are one half, one quarter, and one eighth.
The rebuilt-image test recalculates the brightness boundary.
The fixed-image test retains the baseline boundary.
Column definitions and view selections follow
Table~\ref{tab:imagingvariants}.}
\label{tab:thinning}
\small
\setlength{\tabcolsep}{5pt}
\begin{tabular}{llrrrrrr}
\toprule
Change & Band & $N$ & Median $\Delta b_{\mathrm{link}}$ & $P_{16}$--$P_{84}$
& Median $\Delta\mathrm{IoU}$ & Edge & Below $0.90$ \\
\midrule
Rebuilt image, factor 2 & $V$ & 71 & $0$ & $[-0.02,0]$ & $-0.0057$ & 3 & 24 \\
 & $g$ & 71 & $0$ & $[-0.02,0.016]$ & $-0.0045$ & 4 & 18 \\
Rebuilt image, factor 4 & $V$ & 71 & $0$ & $[-0.02,0.016]$ & $-0.0138$ & 2 & 40 \\
 & $g$ & 68 & $0$ & $[-0.02,0.02]$ & $-0.0108$ & 4 & 34 \\
Rebuilt image, factor 8 & $V$ & 66 & $0$ & $[-0.03,0.02]$ & $-0.0198$ & 6 & 57 \\
 & $g$ & 67 & $0$ & $[-0.0244,0.02]$ & $-0.0157$ & 5 & 47 \\
Fixed image, factor 2 & $V$ & 66 & $0.02$ & $[0.02,0.04]$ & $0.0001$ & 10 & 0 \\
 & $g$ & 66 & $0.02$ & $[0.02,0.03]$ & $0.0002$ & 13 & 1 \\
Fixed image, factor 4 & $V$ & 50 & $0.05$ & $[0.03,0.0616]$ & $0.0002$ & 39 & 0 \\
 & $g$ & 51 & $0.05$ & $[0.03,0.06]$ & $0.0005$ & 40 & 0 \\
Fixed image, factor 8 & $V$ & 30 & $0.07$ & $[0.05,0.0736]$ & $0.0008$ & 96 & 1 \\
 & $g$ & 30 & $0.07$ & $[0.0628,0.08]$ & $0.0007$ & 99 & 0 \\
\bottomrule
\end{tabular}
\end{table*}

The calibration applies to the image construction and particle sampling
specified above.
Tables~\ref{tab:imagingvariants} and \ref{tab:thinning} summarize the
completed sensitivity tests on the original nine-node grid.
Each halo contributes the median selected parameter and median overlap
over three views.
Changes are measured relative to the baseline, retaining only haloes
with interior maxima in both calculations.
The last two columns instead count individual views with an interior
baseline maximum and overlap at least $0.90$, comprising 199 in $V$
and 198 in $g$ across all 76 haloes.
The edge and low-overlap categories are disjoint.
The supplementary tables retain individual curves and exceptional cases.



Removing the PSF wing leaves the halo parameter unchanged in $V$ and
changes only two selected parameters in $g$, each by $-0.02$.
A sky residual of $29\,\mathrm{mag\,arcsec}^{-2}$ corresponds to
$0.16$ per cent of a $22\,\mathrm{mag\,arcsec}^{-2}$ sky and gives
zero median change in the linking parameter.
A residual of $27\,\mathrm{mag\,arcsec}^{-2}$ corresponds to $1$ per
cent and fills every $g$ stamp because the residual equals the threshold.
The resulting endpoint cannot define a calibrated linking parameter.
The 28 $V$ views reaching an endpoint include the qualifying view of
the excluded halo, leaving 27 among the 75 fully interior haloes.

Moving the brightness threshold by half a magnitude changes the selected
parameter and can replace a connected region with a different component.
Consequently, a brighter threshold does not always shrink the selected
region monotonically.
Including satellites represented by at least $100$ stars changes only
the image, with satellite particles still excluded from the link.
The median parameter change is zero, but some selected regions no longer
surround the central galaxy.
The satellite edge counts include the excluded halo and become 0 in $V$
and 1 in $g$ after omitting the halo.



Particle thinning retains a nested random subset while preserving each
retained particle's mass and luminosity.
Rebuilding the image lowers surface brightness and leaves the median
parameter change at zero, despite decreasing overlap and losing detectable
regions.
At a factor of eight, thirteen views in each band have no bright region.
Holding the image fixed instead raises the median parameter by $0.02$,
$0.05$, and $0.07$ at factors of two, four, and eight.
The increasing endpoint counts restrict the last comparison to 30 haloes
per band.
Neither test establishes convergence of the hydrodynamical simulation.

For equal-mass particles, a thinning factor $f$ increases the mean
particle separation by $f^{1/3}$.
Because retained particle masses remain unchanged, preserving a similar
linking distance relative to the particle separation would require
$b_{\mathrm{link}}$ to increase by the same factor.
Starting from $b_{\mathrm{link}}=0.10$, the expected increases are approximately
$0.026$, $0.059$, and $0.10$.
The measured changes have the same sign and order, with a smaller change
at a factor of eight.
A fixed dimensional linking distance would instead leave $b_{\mathrm{link}}$ unchanged
under the implemented mass convention in main-paper Equation~\eqref{main-eq:linkcode}.
The conversion $b_{\mathrm{link}}=0.10$ to $1.046\,h^{-1}\mathrm{ckpc}$ for a
$10^{5}\,h^{-1}\mathrm{M}_{\odot}$ particle is a reference value, rather than a
sample median or a universal density boundary.

In the archived-normalization analysis, the image centre lies below the
brightness threshold in 34 $V$ projections and 35 $g$ projections.
Without the nearest-pixel selection, the regions are empty in those
projections, including 26 in $V$ and 25 in $g$ with baseline overlap at
least $0.90$. The calibration therefore depends on this selection rule.
Dust, additional viewing directions, other simulation volumes, and
fixed $100\,\mathrm{kpc}$ apertures have not been tested. The fixed
30 pkpc aperture is evaluated in the CAMELS transfer test.


\section{Archived additional image diagnostics}
The exploratory diagnostics use the original photometric normalization
and do not replace the corrected-distance primary measurements.
Across all 76 haloes, the median positional light fractions inside the
boundary are $0.7947$ in $V$ and $0.8240$ in $g$.
The corresponding fractions from convolved images are $0.7864$ and
$0.8121$.
The fractions use satellite-excluded total light and differ from completeness
and precision.
The PSF-weighted overlap in main-paper Section~\ref{main-sec:adopted} quantifies the consequence
for the selected linking length.

The equivalent radius is the radius of a circle with the selected area.
Median halo radii are $4.05\,\mathrm{arcsec}$ in $V$ and
$3.96\,\mathrm{arcsec}$ in $g$, with 16th to 84th percentiles
$2.21$--$16.21\,\mathrm{arcsec}$ and $2.34$--$14.95\,\mathrm{arcsec}$.
For the 29 $V$ haloes and 28 $g$ haloes above
$5.0\,\mathrm{arcsec}$, the exterior fractions within
$1.0\,\mathrm{arcsec}$ fall to $0.4342$ and $0.4252$.
The interior fractions fall to $0.1930$ and $0.2120$.
The residual therefore extends beyond the immediate boundary more
strongly in larger regions.

\section{Archived total and radial stellar luminosity}
The numerical summaries use the archived normalization. The main text
reports the corrected primary luminosity comparisons.
We compare the total group light with light from stars projected
inside the isophote using the signed fractional difference
$L_{\mathrm{group}}/L_{\mathrm{inside}}-1$.
For each halo, we take the median over three views before calculating
the median and percentiles across 75 haloes.
At $b_{\mathrm{link}}=0.20$, the median excess is $7.04$ per cent in $V$ and
$6.33$ per cent in $g$. At $b_{\mathrm{link}}=0.10$, the median differences are
$0.31$ and $0.28$ per cent, with 16th to 84th percentiles
$-2.47$ to $4.18$ per cent and $-2.78$ to $4.16$ per cent.
For the absolute differences quoted in the main paper, the absolute
value is taken separately in every view before either median.

The radial table divides each image into projected radii below
$0.5R_{\mathrm{eq}}$, $0.5$--$1.0R_{\mathrm{eq}}$,
$1.0$--$1.5R_{\mathrm{eq}}$, and above $1.5R_{\mathrm{eq}}$,
where $R_{\mathrm{eq}}$ is the radius of a circle with the isophotal
area. The table records group light, light from stars inside the
isophote, and group light inside the isophote for each annulus.
Exterior group light at radii above $R_{\mathrm{eq}}$, divided by
$L_{\mathrm{inside}}$, has medians $6.60$ and $5.86$ per cent at
$b_{\mathrm{link}}=0.20$, falling to $1.38$ and $1.47$ per cent at $b_{\mathrm{link}}=0.10$ in $V$
and $g$, respectively. An equivalent radius approximates the area
and does not reproduce the shape of a non-circular isophote.
All luminosity sums use full stellar fluxes selected by projected
position. The sums do not integrate PSF-convolved aperture flux.

\section{Overlap near the isophote}
Let $G$ denote the stellar group, $I$ the stars projected inside the
isophote, and $R$ the stars projected into the boundary band.
The band contains one layer of pixels on each side of the boundary,
including pixels touching across a corner. The overlap in the band is
defined as
\begin{equation}
\mathrm{IoU}_{R}=\frac{L(G\cap I\cap R)}{L[(G\cup I)\cap R]}.
\end{equation}
Completeness divides $L(G\cap I\cap R)$ by $L(I\cap R)$, and precision
divides the same numerator by $L(G\cap R)$. Every sum uses full stellar
fluxes selected by projected position. A ratio with a zero denominator
is undefined. The halo statistic takes the median over all three views
and remains undefined when any view has an undefined ratio.
The pixel band spans different physical widths at different redshifts.

At the parameter selected from the corrected-distance total light, the
median overlaps in the band are $0.4614$ in $V$ and $0.4760$ in $g$.
The 16th to 84th percentiles are $0.3924$--$0.5201$ and
$0.3873$--$0.5399$. At the selected $b_{\mathrm{link}}$, completeness is $0.6945$
($0.5789$--$0.7949$) in $V$ and $0.7004$ ($0.5603$--$0.8018$) in $g$.
Precision is $0.5816$ ($0.5301$--$0.6528$) and $0.5864$
($0.5351$--$0.6586$), respectively. Parentheses give the 16th--84th
percentile ranges across halo medians. The corrected photometry changes
some projection-level selected parameters and therefore changes the light ratios.
All 75 halo medians are defined in both bands.
Table~\ref{tab:ring} also gives completeness and precision.
The larger groups at $b_{\mathrm{link}}=0.20$ recover more light from stars just
inside the isophote, increasing completeness in the band.
The lower precision near the boundary partly offsets the gain.

\begin{table}
\centering
\caption{Median light ratios in the pixel band on both sides of the
isophote. $C$ denotes completeness and $P$ precision.
Selected $b_{\mathrm{link}}$ maximizes the ratio using all stellar light separately
in each view. Each entry first takes the median over three views,
then over the 75 retained halo entries. All ratios are defined.}
\label{tab:ring}
\small
\setlength{\tabcolsep}{2.5pt}
\begin{tabular}{lrrrrrr}
\toprule
Choice & $\mathrm{IoU}_{V}$ & $C_{V}$ & $P_{V}$ & $\mathrm{IoU}_{g}$ & $C_{g}$ & $P_{g}$ \\
\midrule
Selected $b_{\mathrm{link}}$ & 0.4614 & 0.6945 & 0.5816 & 0.4760 & 0.7004 & 0.5864 \\
$b_{\mathrm{link}}=0.10$ & 0.4936 & 0.7648 & 0.5651 & 0.4753 & 0.7509 & 0.5697 \\
$b_{\mathrm{link}}=0.20$ & 0.5410 & 0.9703 & 0.5470 & 0.5480 & 0.9722 & 0.5572 \\
\bottomrule
\end{tabular}
\end{table}

\subsection{Sensitivity to boundary-band width}
\label{supp:boundary-width}

For each image pixel, we calculate the signed Euclidean distance between
its centre and the centre of the nearest pixel in the opposite isophotal
class. Pixels with $|d|\leq\Delta$ form the boundary band. Thus $\Delta$
is a pixel-centre distance threshold, not an exact distance from a
continuous contour. We repeat the positional-flux overlap measurement at
thresholds $\Delta=0.2$, $0.3$, $0.5$, $1$, and $2$ arcsec. The $0.5$-arcsec
threshold is comparable to the $1$-arcsec PSF FWHM, but it does not define
an exact one-FWHM-wide band. For each view, the position-calibrated value
is selected using the original nine-point full-image positional grid; we
compare it with fixed $b_{\mathrm{link}}=0.20$. Both the overlap and which
parameter gives the higher median depend on the distance threshold. The
fixed value $0.20$ gives higher median overlap at thresholds up to $1$ arcsec
in $V$, and up to $0.5$ arcsec in $g$. At a $2$-arcsec threshold, the
position-calibrated value gives higher median overlap in both bands. The shaded regions in
Fig.~\ref{fig:supp-boundary-width} show the 16th--84th percentiles across
halo-entry medians.

\begin{table*}
\centering
\caption{Median light-weighted overlap within boundary bands defined by
pixel-centre distance thresholds. $\Delta$ is the maximum distance from a
pixel centre to the nearest pixel centre of the opposite isophotal class;
it is not an exact distance from a continuous contour. The selected
parameter is the full-image positional optimum on the original nine-point
grid for each view; the comparison uses fixed $b_{\mathrm{link}}=0.20$.
The last column is the difference between the two medians.}
\label{tab:supp-boundary-width}
\small
\begin{tabular}{lrrrr}
\toprule
Band & Distance threshold $\Delta$ (arcsec) & Selected $b_{\mathrm{link}}$ &
$b_{\mathrm{link}}=0.20$ & Difference \\
\midrule
$V$ & 0.2 & 0.441 & 0.526 & $-0.085$ \\
$V$ & 0.3 & 0.461 & 0.541 & $-0.080$ \\
$V$ & 0.5 & 0.516 & 0.575 & $-0.059$ \\
$V$ & 1.0 & 0.651 & 0.671 & $-0.021$ \\
$V$ & 2.0 & 0.842 & 0.808 & $+0.035$ \\
$g$ & 0.2 & 0.453 & 0.531 & $-0.078$ \\
$g$ & 0.3 & 0.476 & 0.548 & $-0.072$ \\
$g$ & 0.5 & 0.528 & 0.592 & $-0.065$ \\
$g$ & 1.0 & 0.699 & 0.696 & $+0.003$ \\
$g$ & 2.0 & 0.862 & 0.851 & $+0.012$ \\
\bottomrule
\end{tabular}
\end{table*}

\begin{figure*}
\centering
\includegraphics[width=\textwidth]{supplement/astra_review_response_20261106/results/boundary_width_sensitivity.pdf}
\caption{Median overlap in bands around the isophote as a function of
pixel-centre distance threshold $\Delta$. Blue points use the full-image
positional optimum on the original nine-point grid for each halo view;
orange points use $b_{\mathrm{link}}=0.20$. Shading shows the 16th--84th
percentiles across the 75 halo-entry medians. The dashed line marks
$\Delta=0.5$ arcsec, a threshold comparable to the $1$-arcsec PSF FWHM.
The band is defined by distance between pixel centres, not by a continuous
contour distance.}
\label{fig:supp-boundary-width}
\end{figure*}

\section{Photometric provenance}
The archived grids contain no record of the generating software version.
We reproduce all $32\times12\times2$ flux values in each of the seven
grids with python-fsps 0.4.7 and the recorded population parameters.
Every reproduced flux equals the archived value at stored precision.
The reproducing installation reports the \texttt{mist} isochrone,
\texttt{miles} spectral, and \texttt{DL07} dust-emission libraries.
The attenuation parameters \texttt{dust1} and \texttt{dust2} are zero.
The default circumstellar dust model for asymptotic-giant-branch stars
remains enabled. The omitted attenuation concerns interstellar dust.
The supplementary metadata record the compiled extension's hash.
The available FSPS source uses $\Omega_m=0.27$,
$\Omega_\Lambda=0.73$, and $H_0=72\,\mathrm{km\,s}^{-1}\mathrm{Mpc}^{-1}$
for luminosity distances. Angular sizes use the simulation cosmology.
The archived image calculation applies no additional distance factor.

We recover the FSPS distance table using the source integration and
redshift interpolation. A consistent distance convention multiplies
each apparent stellar flux by
$[D_{L,\mathrm{FSPS}}(z)/D_{L,\mathrm{sim}}(z)]^2$.
The flux multipliers are $0.8873$, $0.8895$, $0.9041$, $0.9249$,
$0.9476$, and $0.9589$ at the outputs from lowest to highest redshift.
The corresponding magnitude shifts are $0.1298$, $0.1272$, $0.1095$,
$0.0848$, $0.0584$, and $0.0456$.
The correction preserves stellar spectra, angular scales, the PSF,
and the numerical brightness limits. The common flux factor cancels
from luminosity ratios for fixed stellar selections but changes the
isophotal regions selected from the images.
The absolute-magnitude grid used for the redshift-zero control needs
no luminosity-distance correction.

Recalculating all 76 haloes with the common distance convention
retains the same 75 halo entries in both bands. The median parameter
changes for 8 entries in $V$ and 11 in $g$. The median change per
entry is zero in both bands, including the 16th and 84th percentiles.
The sample median changes from $0.10$ to $0.08$ in $V$ and remains
$0.10$ in $g$. A few changes can move the central entry of the ordered
sample across two adjacent grid values without moving the median of
the paired differences.

Table~\ref{tab:distance} gives the parameters at each output.
The median overlaps become $0.9658$ in $V$ and $0.9721$ in $g$, with
$b_{\mathrm{link}}$ chosen separately for each view. The median paired changes in
overlap are $-0.0012$ and $-0.0009$.
The median ratios of corrected to original region area are $0.9386$
and $0.9425$ across all 228 views per band. No view loses every pixel
brighter than the limit.
The area increases in one view per band because the nearest-pixel
rule selects a different region. The $g$ case is the single-pixel
example in main-paper Section~\ref{main-sec:fail}.
In $V$, the excluded halo at $z=0.227$ changes from $0.04$ to
$0.88\,\mathrm{arcsec}^{2}$ along $x$. The overlap remains below
$10^{-5}$ and the selected $b_{\mathrm{link}}$ remains at the upper grid endpoint.

\begin{table}
\centering
\caption{Median linking parameters after using the simulation cosmology
for both angular and luminosity distances. $N$ counts the same halo
entries retained in the original calibration. Each halo value first
takes the median of the three parameters selected separately by view.}
\label{tab:distance}
\begin{tabular}{lrrr}
\toprule
$z$ & $N$ & $b_{V}$ & $b_{g}$ \\
\midrule
3.002 & 1 & 0.10 & 0.12 \\
1.990 & 3 & 0.10 & 0.15 \\
1.002 & 17 & 0.08 & 0.08 \\
0.500 & 33 & 0.08 & 0.08 \\
0.227 & 9 & 0.12 & 0.10 \\
0.190 & 12 & 0.12 & 0.10 \\
\bottomrule
\end{tabular}
\end{table}



\section{Archived-normalization illustrations retained for reference}
\label{sec:archived-products}
\setcounter{figure}{0}
\setcounter{table}{0}
\renewcommand{\thefigure}{S\arabic{figure}}
\renewcommand{\thetable}{S\arabic{table}}

The following figures and table preserve earlier illustrative products
made with the archived FSPS luminosity-distance normalization. The
corrected-distance calculations in the main paper supersede the numerical
calibration inferred from those images. The archival products do not enter
the primary measurements or conclusions.

\subsection{Representative cluster image}
The image shows how the stellar group changes as the linking parameter
increases. Its contour and displayed overlap use the archived
normalization; the corrected calibration is reported in the main text.
\begin{figure*}
\centering
\includegraphics[width=\textwidth]{figures/cluster_38857678.pdf}
\caption{The most massive halo in the sample at $z=0.190$, with group
mass $1.337\times10^{14}\,h^{-1}\mathrm{M}_{\odot}$ and $38857678$ stellar
particles, projected along the box $y$ axis.
Panels (a)--(e) sample stars with probability proportional to stellar
flux after removing satellite galaxies represented by at least 100
stars. Colours distinguish group members from the remaining stars.
The fixed contour bounds the region brighter than
$27\,\mathrm{mag\,arcsec}^{-2}$ in $g$, with area
$2896.240\,\mathrm{arcsec}^{2}$.
Maximum overlap in the illustrated view occurs at $b_{\mathrm{link}}=0.10$. The image,
contour, and area use the archived photometric normalization. With the corrected distance convention, the median selected
$b_{\mathrm{link}}$ is $0.12$ in $V$ and $0.10$ in $g$.
Panel (f) includes satellite stars and shows stellar mass on a
logarithmic colour scale. Ellipses show second-moment shapes for 935 galaxies,
with semi-axes enlarged by a factor of five.
The dashed square locates panels (a)--(e). Scale bars give proper distances.}
% Alt text: Six views of the most massive simulated group show stellar light and group membership at progressively larger linking parameters, alongside the full stellar-mass map and galaxy shapes. A fixed surface-brightness contour marks the apparent central-galaxy boundary.
\label{supp:cluster-archived}
\end{figure*}

\subsection{Image-region failure examples}
The maps illustrate small or displaced regions selected by the original
centre-based rule. Their panel annotations show archived overlap measurements. The main
text gives the corresponding corrected-distance result for the one-pixel
example.
\begin{figure*}
\centering
\includegraphics[width=\textwidth]{figures/fail_maps.pdf}
\caption{Archived-normalization illustrations of isophotes enclosing
small areas or regions far from the stellar density seed.
Scale bars give angular distances in arcseconds in every panel.
Labels inside each panel give the band, viewing direction, linking
parameter, and overlap.
The common grey scale gives convolved surface brightness.
Cyan plus signs mark image centres, orange crosses mark visible density
seeds, and open purple circles mark watershed cores.
An arrow in panel (a) points towards a core outside the window.
Blue points show stellar group members. Teal contours enclose regions
containing the density seed, and coral contours enclose regions chosen
by the original image-centre rule.
Panel (a) shows the one-pixel $g$ region along $y$ at $z=1.002$.
Panels (b) and (c) show the excluded halo at $z=0.227$ in $V$ along
$x$ and in $g$ along $z$, enlarged to $\pm60\,\mathrm{arcsec}$.
Panels (d)--(f) show the low-overlap halo at $z=1.990$ in $V$ along
$x$ and $y$ and in $g$ along $z$.
The panel labels show archived overlap measurements. The text also reports the
corrected-distance change for the one-pixel example.}
% Alt text: Six image cutouts illustrate tiny or displaced surface-brightness regions at high redshift. Markers identify the image centre, stellar-density seed, watershed cores, and selected group stars; paired contours compare the original region choice with the region containing the seed.
\label{supp:fail-archived}
\end{figure*}

\subsection{Archived boundary-distance fractions}
The fractions use archived fluxes and the archived view-specific linking
parameters. Ring includes pixels touching either side of the contour,
including corner contacts. For each halo, the table reports the median
over three views.
\begin{table}
\centering
\caption{Archived-normalization light near the isophote from group members outside the boundary
and non-members inside. Entries are median fractions across 75 haloes
with maximum overlap between the grid endpoints in all three views.
Out fractions divide group light within the stated distance outside
the isophote by all group light outside. In fractions divide light
from non-members within the stated distance inside the isophote by all
non-member light inside. Ring uses pixels touching the boundary on
either side, including corner contacts. The supplementary table gives
the fluxes for individual haloes.}
\label{supp:boundarydistance}
\small
\begin{tabular}{lrrrr}
\toprule
Selection & Out $V$ & Out $g$ & In $V$ & In $g$ \\
\midrule
Ring & 0.2764 & 0.2632 & 0.1400 & 0.1441 \\
Within $1.0''$ & 0.7300 & 0.7012 & 0.5242 & 0.5501 \\
Within $2.0''$ & 0.9468 & 0.9365 & 0.8318 & 0.8671 \\
Within $5.0''$ & 1 & 1 & 1 & 1 \\
\bottomrule
\end{tabular}
\end{table}

\subsection{No-dimming and unredshifted-spectrum control}
\label{supp:rest-control}

We make images without cosmological dimming and with unredshifted
spectra for the halo at $z=3.002$ and the three haloes at $z=1.990$.
We keep the angular pixel scale, PSF, and numerical brightness limits.
The stellar populations are evaluated at redshift zero, and absolute
AB fluxes per proper square parsec are converted to surface brightness
with $2.5\log_{10}[(206265/10)^{2}]$.
The comparison changes both dimming and sampled wavelength.
For example, the observed $V$ filter samples wavelengths near
$1400\,\text{\AA}$ at $z=3.002$.
The archived intrinsic-surface-brightness images differ from the
artificial-redshift \texttt{rest\_sed\_no\_dimming} stress test. The stress
test places absolute 10-pc stellar flux directly into angular pixels
without conversion to flux per proper area.

We extend the grid to $b_{\mathrm{link}}=0.25$, $0.30$, $0.40$, and $0.50$.
In all 24 combinations of band and view, overlap continues to increase
at the largest tested parameter, $b_{\mathrm{link}}=0.50$. The curves therefore do not
reach a maximum within the tested range.
Precision at the endpoint is $0.9931$ to $1.0000$.
Nearly all group light comes from stars inside the isophote, but some
stars inside the isophote remain outside the group.
Figure~\ref{supp:rest-archived} shows the archived-normalization control curves.
The exploratory curves do not enter the corrected primary calibration.
Image construction therefore changes the overlap curves.
The small sample cannot establish a trend with redshift.

\begin{figure}
\centering
\includegraphics[width=\columnwidth]{figures/b_rest.pdf}
\caption{Overlap without cosmological dimming and with unredshifted
spectra at the same numerical brightness limits. Lines are medians over
three views, and shaded bands span the view range. Colour denotes
redshift, while line style distinguishes haloes at $z=1.990$.
The legend gives group masses in $h^{-1}\mathrm{M}_{\odot}$.
All view curves continue to rise at the final grid point, $b_{\mathrm{link}}=0.50$.}
% Alt text: Overlap curves for four high-redshift haloes rise through the largest tested linking parameter, 0.50. Colours distinguish redshift and line styles distinguish the haloes at redshift 1.990; the shaded ranges show variation across three views.
\label{supp:rest-archived}
\end{figure}


\FloatBarrier
\begin{thebibliography}{99}
\bibitem[Kim et al.(2023)]{kim2023}
Kim, J., Lee, J., Laigle, C., et al. 2023, ApJ, 951, 137
\bibitem[Marinacci et al.(2018)]{marinacci2018tng}
Marinacci, F., et al. 2018, MNRAS, 480, 5113
\bibitem[More et al.(2011)]{more2011}
More, S., Kravtsov, A.~V., Dalal, N., \& Gottl\"ober, S. 2011, ApJS, 195, 4
\bibitem[Naiman et al.(2018)]{naiman2018tng}
Naiman, J.~P., et al. 2018, MNRAS, 477, 1206
\bibitem[Nelson et al.(2018)]{nelson2018tng}
Nelson, D., et al. 2018, MNRAS, 475, 624
\bibitem[Nelson et al.(2019)]{nelson2019tng}
Nelson, D., et al. 2019, Comput.\ Astrophys.\ Cosmol., 6, 2
\bibitem[Pillepich et al.(2018)]{pillepich2018tng}
Pillepich, A., et al. 2018, MNRAS, 475, 648
\bibitem[Rodriguez-Gomez et al.(2015)]{rodriguezgomez2015}
Rodriguez-Gomez, V., Genel, S., Vogelsberger, M., et al. 2015, MNRAS, 449, 49
\bibitem[Springel et al.(2018)]{springel2018tng}
Springel, V., et al. 2018, MNRAS, 475, 676
\end{thebibliography}

\bsp
\end{document}
